\documentclass[12pt, a4paper, openany]{book}

\usepackage{fontspec}
\setmainfont{TeX Gyre Pagella}
\usepackage{geometry}
\geometry{margin=1.25in, top=1.2in, bottom=1.3in}
\usepackage{setspace}
\onehalfspacing
\usepackage{amsmath}
\usepackage{amssymb}
\usepackage{amsthm}
\usepackage{mathtools}
\usepackage{stmaryrd}
\usepackage{titlesec}
\usepackage{titletoc}
\usepackage{fancyhdr}
\setlength{\headheight}{14pt}
\usepackage{booktabs}
\usepackage{array}
\usepackage{longtable}
\usepackage{listings}
\usepackage{xcolor}
\usepackage{csquotes}
\usepackage{enumitem}
\usepackage{hyperref}
\usepackage{parskip}
\usepackage{microtype}

\hypersetup{colorlinks=true, linkcolor=black, citecolor=black, urlcolor=black}

\setlength{\parindent}{1.5em}
\setlength{\parskip}{0pt}

% Theorem environments
\theoremstyle{plain}
\newtheorem{theorem}{Theorem}[chapter]
\newtheorem{proposition}[theorem]{Proposition}
\newtheorem{lemma}[theorem]{Lemma}
\newtheorem{corollary}[theorem]{Corollary}
\theoremstyle{definition}
\newtheorem{definition}[theorem]{Definition}
\newtheorem{example}[theorem]{Example}
\newtheorem{construction}[theorem]{Construction}
\theoremstyle{remark}
\newtheorem{remark}[theorem]{Remark}

% Chapter and section formatting
\titleformat{\chapter}[display]
  {\normalfont\large\bfseries}
  {\chaptertitlename\ \thechapter}{12pt}{\Large}
\titleformat{\section}{\normalfont\large\bfseries}{}{0em}{}
\titleformat{\subsection}{\normalfont\normalsize\bfseries}{}{0em}{}
\titleformat{\subsubsection}{\normalfont\normalsize\itshape}{}{0em}{}

% Code listings
\definecolor{codegreen}{rgb}{0,0.5,0}
\definecolor{codegray}{rgb}{0.5,0.5,0.5}
\definecolor{codeblue}{rgb}{0,0,0.6}
\lstdefinestyle{rust}{
  language=C,
  basicstyle=\small\ttfamily,
  keywordstyle=\color{codeblue}\bfseries,
  commentstyle=\color{codegray}\itshape,
  stringstyle=\color{codegreen},
  breaklines=true,
  frame=single,
  framesep=4pt,
  xleftmargin=12pt,
  xrightmargin=4pt,
}

% Operators
\newcommand{\Pop}{\mathsf{Pop}}
\newcommand{\Refuse}{\mathsf{Refuse}}
\newcommand{\Bind}{\mathsf{Bind}}
\newcommand{\Collapse}{\mathsf{Collapse}}
\newcommand{\Hist}{\mathbf{Hist}}
\newcommand{\Obs}{\mathbf{Obs}}
\newcommand{\calH}{\mathcal{H}}
\newcommand{\calC}{\mathcal{C}}
\newcommand{\calA}{\mathcal{A}}
\newcommand{\calN}{\mathcal{N}}
\newcommand{\calR}{\mathcal{R}}
\newcommand{\Adm}{\mathsf{Admissible}}
\newcommand{\Rfsd}{\mathsf{Refused}}
\newcommand{\Col}{\mathsf{Collapsed}}
\newcommand{\Proc}{\mathsf{Process}}

\pagestyle{fancy}
\fancyhf{}
\fancyhead[LE]{\small\textit{Spherepop: Toward a Complete Theory}}
\fancyhead[RO]{\small\textit{\leftmark}}
\fancyfoot[C]{\thepage}

\title{
  \vspace{2cm}
  {\Huge\bfseries Spherepop}\\[1em]
  {\large\itshape Histories, Continuations, and the Algebra\\
  of Reachable Computation}\\[3em]
  {\normalsize Toward a Complete Theory: History-Primary Computation,\\
  Type Systems, Error Correction, and Garbage Collection}
}
\author{Flyxion\\[0.5em]\small Independent Researcher}
\date{}


\begin{document}

\maketitle
\thispagestyle{empty}

\frontmatter

\chapter*{Abstract}
\addcontentsline{toc}{chapter}{Abstract}

Most programming languages begin with values, variables, and state transitions.
Spherepop begins elsewhere: with the claim that histories are ontologically prior
to states, and that computation is the progressive restriction of possibility rather
than the manipulation of values.

This monograph presents the complete theory of the Spherepop programming language
and its philosophical foundations. We proceed from two converging arguments. The
first is a theory of notes: the observation that notes are not passive records of
the past but anti-collapse artifacts that preserve access to continuations. The
second is the Spherepop inversion: the observation that state is a collapsed
history and that taking this seriously demands a new computational ontology.

These arguments converge because they are the same argument. A note is an
externalized continuation. A Spherepop history is an internalized note. Civilization
builds note infrastructure; Spherepop builds note infrastructure for computation.

The monograph develops this convergence through nine parts. Part~I establishes
the dual failure of the passive theory of notes and the state-centric theory of
computation. Part~II derives the four Spherepop operators---Pop, Refuse, Bind, and
Collapse---from first principles as the minimal complete set for managing possibility
space. Part~III develops the full operational semantics, the History category, and
collapse as quotient. Part~IV presents the complete admissibility type system with
Progress, Preservation, and Collapse Soundness theorems. Part~V derives the
seven note classes as physical realizations of the four operators, establishing
that notes are Spherepop programs and Spherepop programs are notes. Part~VI covers
compilation, the Event IR, history equivalence as compiler correctness, garbage
collection as admissibility-guided collapse, and error correction as certified
refusal. Part~VII develops the deep mathematical structure: sheaf semantics, CLIO
projections, the Santoro--Waghmare--Panaretos unification, and active geodesic
inference. Part~VIII situates the framework within civilization-scale note
infrastructure, MEM\textbar8, and Repair Theory. Part~IX collects open problems.

Four principal theorems anchor the work: the Conservation of Possibility, the
Ephemerality Inversion, the Civilizational Collapse Theorem, and the
History--Observation Adjunction. Four implementation discoveries retroactively
sharpened the theory: refusal must carry reasons; collapse must specify quotient
rules; type checking depends on the assumed boundary of the world; and compiler
correctness requires history equivalence, not mere output equivalence.


\begin{quote}\small\textit{This manuscript is not presented as a finished formalization. It is a staged research program. Some results are proved fully, some are proof sketches, and some are conjectural bridges to be developed through implementation and subsequent work. The goal is to make the program precise enough to be falsifiable, not to claim it complete.}\end{quote}

\newpage

\tableofcontents

\mainmatter

%% ============================================================
%% PART I: THE PROBLEM
%% ============================================================
\part{The Problem}

\chapter{The Passive Theory of Notes and Its Failure}

There is a theory of notes so pervasive that it rarely receives explicit statement.
On this theory, a note is a device for compensating the limitations of biological
memory. Something happens---a meeting is scheduled, a formula is derived, a
perception is formed---and because the mind cannot reliably retain it, the note
is made. The note is therefore a prosthetic memory: it externalizes information
onto a more durable substrate, where it may be retrieved when biological recall
fails.

This is the passive theory of notes. It places the note in a retrospective
relation to time. The note exists because something happened in the past that is
worth preserving. Its value is proportional to the fidelity with which it
reproduces a prior state: a good note accurately records what occurred; a poor note
distorts it; an ephemeral note is low-value because it preserves little and does
not endure.

The passive theory is not entirely wrong. Many notes are made for the reason it
describes. But as a general account of what notes are and why they matter, it
fails on its own terms and in a direction that reveals something important about
cognition, language, and civilization.

The failure becomes apparent the moment one examines not celebrated notes but
common ones. A grocery list does not record anything that happened. The items on
the list do not exist yet in the form the list anticipates: they have not been
purchased, assembled, or carried home. The list records a future act. A blueprint
describes a building not yet constructed. Source code records computations not yet
run. A theorem records consequences not yet derived. A calendar records commitments
not yet honored. A playlist records a future act of listening. A bookmark records
a future act of reading.

None of these are memories. All of them are notes.

The passive theory cannot accommodate this class of cases without becoming
incoherent. It might be extended to say that a note records an \emph{intention}
rather than an event, but this extension merely postpones the difficulty. The
deeper problem is that the passive theory places every note in a backward-looking
relation to time. Even when the recorded content is a future intention, the note
is still conceived as a trace of a prior mental event. The orientation is always
retrospective.

This gets the phenomenology of note-taking almost exactly wrong. When one makes a
to-do list, one is not primarily recording a past intention; one is constructing
a future. The note is not oriented toward what was; it is oriented toward what
has not yet occurred and what might not occur without the note's assistance.

\section{The Continuation Account}

A more adequate theory begins from a distinction that theoretical computer science
has found indispensable: the difference between a \emph{state} and a
\emph{continuation}. A state is a snapshot of a system at a moment in time---a
summary, a compression, a projection. A continuation is what comes next: the
future execution enabled by a present configuration, the program of further
development that a state makes possible.

Most theories of information focus on states. A database stores states. A
photograph is described as capturing a state. A memory is treated as the retention
of a past state. But what makes any stored state useful is not the state itself.
It is the continuation that state makes accessible.

Consider a Git commit. The commit is not valuable because it stores bytes. It is
valuable because it preserves a path through development space---a specific
trajectory through the possibility space of a codebase that would otherwise require
expensive reconstruction. The stored state is instrumentally useful only because
it is a node in a navigable graph of continuations.

\begin{definition}[Note]
Let $\calC$ denote a space of continuations and let $A$ be an agent. A
\emph{note} is any artifact $N$ such that there exists a continuation
$c \in \calC$ for which
\[
P_A(c, t \mid N) > P_A(c, t),
\]
where $P_A(c, t)$ denotes the probability that agent $A$ can successfully
reconstruct or traverse continuation $c$ at future time $t$.
\end{definition}

This definition is deliberately broad. It immediately encompasses bookmarks, maps,
photographs, source code, language, and libraries without privileging any medium.
What unifies them is the functional condition: a note increases the reachability
of at least one continuation.

\begin{definition}[Continuation Volume]
The \emph{continuation volume} of a note $N$ is
\[
V_C(N) = \bigl|\{c \in \calC : P_A(c, t \mid N) > P_A(c, t)\}\bigr|.
\]
\end{definition}

\begin{definition}[Reachability Leverage]
The \emph{reachability leverage} of a note $N$ is
\[
\Lambda(N) = \frac{\sum_{c} \bigl[C_r(c) - C_r(c \mid N)\bigr]}{C_{\text{create}}(N)},
\]
where $C_r(c)$ is the cost of reconstructing continuation $c$ from first principles
and $C_r(c \mid N)$ is the cost given access to $N$.
\end{definition}

Reachability leverage varies over many orders of magnitude. A ten-second bookmark
may save hours of future search. A one-line equation may preserve centuries of
accumulated derivation. The leverage analysis predicts which notes are most worth
making and preserving, entirely independently of their fidelity to past states.

\section{Notes as Anti-Collapse Artifacts}

The Spherepop framework, developed in subsequent chapters, holds that histories
are primary and states are their compressed residues. A state $s$ is a projection
\[
\sigma : \calH \to S
\]
that collapses many histories $h \in \calH$ into a single observable state,
producing what the framework calls the \emph{state illusion}: the appearance that
the present configuration is self-sufficient, when it is in fact the endpoint of
a history whose fibers have been discarded.

A note is, in this framework, an anti-collapse artifact: an artifact that preserves
selected fibers of $\calH$ against the projection $\sigma$.

\begin{proposition}[Notes as Anti-Collapse Artifacts]
A note $N$ is an anti-collapse artifact if it preserves selected fibers of $\calH$
such that a future agent can partially reconstruct elements of $\calH$ rather than
being restricted to $\sigma(\calH)$.
\end{proposition}

This is the precise sense in which the passive theory inverts the correct account.
The passive theory holds that a note preserves the past against forgetting. The
correct account holds that a note preserves access to futures against collapse.
The distinction is not merely philosophical: it determines what kinds of notes
matter, how notes should be evaluated, and what it means for a note to succeed
or fail.

\chapter{The State-Centric Assumption and Its Consequences}

\section{The Standard Premise}

Open any introductory programming text and you encounter a first lesson: a variable
is a named location that holds a value. A program begins in some initial state,
instructions modify that state, and the final state is the answer.

\begin{lstlisting}[style=rust, caption={The standard state-centric view}]
x = 5
y = x + 1
\end{lstlisting}

This appears natural. It mirrors how we naively think about physical processes: an
object has properties; those properties change; the current properties constitute
the object's state. But the appearance of naturalness is produced by a prior
choice---the choice to treat present configuration as the fundamental unit of
description.

\section{The State Illusion}

Consider what this choice erases. In version control, the repository is not merely
its current file tree; it is the entire sequence of commits that produced it. In
event sourcing, database records are append-only event logs; the current state is
a fold over that log. In scientific experiment, a measurement is the terminus of
a sequence of decisions, calibrations, and interpretive choices. In evolutionary
biology, a genome is a compressed record of selective pressures.

In each case, the historical structure is epistemically and causally prior to the
present configuration. State is the summary. History is the substance.

\begin{definition}[State Degeneracy]
Let $E$ be an event alphabet, $\calH = \bigcup_{n=0}^{\infty} E^n$ the set of all
finite histories, and $S$ an observable state space. A \emph{state projection} is
a function $\sigma : \calH \to S$. The \emph{degeneracy} of a state $s \in S$ is
\[
D(s) = |\sigma^{-1}(s)|.
\]
\end{definition}

\begin{proposition}[The State Illusion]
\label{prop:state-illusion}
For any nontrivial finite state space $S$ and any non-injective projection
$\sigma : \calH \to S$, there exist distinct histories $H_1 \neq H_2$ such that
$\sigma(H_1) = \sigma(H_2)$. Consequently,
\[
\sigma(H_1) = \sigma(H_2) \not\Rightarrow H_1 = H_2.
\]
\end{proposition}

\begin{proof}
Since $|\calH|$ is countably infinite and $|S|$ is finite, $\sigma$ cannot be
injective. Therefore $D(s) > 1$ for at least one $s \in S$, which yields the
required $H_1 \neq H_2$ with $\sigma(H_1) = \sigma(H_2)$.
\end{proof}

\begin{remark}
Proposition~\ref{prop:state-illusion} is the formal statement of the state
illusion: any program that exposes only $\sigma(H)$ to its user has silently
discarded distinguishing information about the history $H$. The conventional
debugger is an instrument for reconstructing $H$ from $\sigma(H)$ after
information has already been lost.
\end{remark}

\section{Consequences for Language Design}

Programming languages that treat state as primary inherit the degeneracy of
Proposition~\ref{prop:state-illusion} as a structural feature. Race conditions
arise because $\sigma$ is non-injective on interleaved histories. Use-after-free
vulnerabilities arise because $\sigma$ conflates histories that differ only in
deallocation events. The auditing and provenance tools that practitioners add to
production systems are, in each case, partial reconstructions of $H$ from $\sigma(H)$.

The Spherepop design decision follows directly: if history is what is lost,
history must be what is preserved.

\section{The Connection to Notes}

The state illusion and the passive theory of notes are the same mistake viewed from
different angles.

The passive theory treats a note as a record of a past state. The state-centric
theory of computation treats a program's output as its definitive product. Both
identify value with a terminal state rather than with the trajectory that produced
it. Both discard the history in favor of the summary.

The Spherepop framework and the continuation theory of notes are therefore the
same correction viewed from different angles. Both insist that the trajectory is
primary. Both insist that what matters is not what a system is at a moment but
what it has done and what it can still do. Both treat state as a derived notion:
useful for some purposes, but never the fundamental unit of description.

This convergence is the organizing principle of the present monograph.

%% ============================================================
%% PART II: FIRST PRINCIPLES
%% ============================================================
\part{First Principles}

\chapter{Deriving the Operators}

\section{The Minimal Requirements}

If histories are primary and states are derived, then a computational framework
must at minimum be able to:

\begin{enumerate}
\item Commit to a specific continuation from among the available possibilities.
\item Document that a possible continuation is inadmissible, with reasons.
\item Record a dependency between two elements of the history.
\item Observe the current state of a history under a specific observational framework.
\end{enumerate}

We claim that these four requirements are not only necessary but sufficient: they
constitute the minimal complete set of operations on possibility space. All other
computational phenomena are derived from these four.

\begin{definition}[Event Alphabet]
Let $E$ be a set of events, partitioned into primitive generators:
\[
E = \{\Pop\} \cup \{\Refuse\} \cup \{\Bind\} \cup \{\Collapse\} \cup E_\lambda
\]
where $E_\lambda$ contains the lambda-calculus events $\{\mathsf{LamIntro},
\mathsf{Apply}\}$ and scope markers.
\end{definition}

\begin{definition}[Spherepop World]
A \emph{Spherepop world} is a pair $W = (H, \Omega)$ where $H \in \calH$ is a
history (finite event sequence) and $\Omega \subseteq \Omega_0$ is the current
option space (a finite set of available symbols), derived from some initial option
space $\Omega_0$ with $|\Omega_0| < \infty$.
\end{definition}

\section{The Four Operators}

\begin{definition}[Pop: Commitment]
The \emph{commitment operator} for symbol $x \in \Omega$ is
\[
P_x : (H, \Omega) \mapsto (H \mathbin{++} [\Pop(x)],\; \Omega \setminus \{x\}).
\]
$P_x$ is undefined when $x \notin \Omega$. Pop records a commitment and removes
the committed symbol from the option space, foreclosing all futures in which $x$
takes a different value.
\end{definition}

\begin{definition}[Refuse: Documented Inadmissibility]
The \emph{refusal operator} for symbol $x$ with reason $r$ is
\[
R_{x,r} : (H, \Omega) \mapsto (H \mathbin{++} [\Refuse(x, r)],\; \Omega).
\]
Note that $\Omega$ is unchanged: refusal records inadmissibility without
foreclosing structural availability. The reason $r$ is a first-class component
of the operation. A refusal without a reason documents nothing.
\end{definition}

\begin{definition}[Bind: Dependency Declaration]
The \emph{bind operator} is
\[
B_{a,b} : (H, \Omega) \mapsto (H \mathbin{++} [\Bind(a, b)],\; \Omega).
\]
Bind records that two trajectory elements are coupled without consuming either
from the option space.
\end{definition}

\begin{definition}[Collapse: Observation under Rule]
The \emph{collapse operator} with collapse rule $c$ is
\[
C_{x,c} : (H, \Omega) \mapsto (H \mathbin{++} [\Collapse(x, c)],\; \Omega),
\]
subject to the admissibility precondition $\Gamma \vdash t : \Adm(T)$.
Collapse requires an admissibility certificate. A term whose type is not
$\Adm(\cdot)$ cannot be collapsed---this is a hard type error, not a runtime
exception.
\end{definition}

\section{Why These Four and No Others}

The minimality of the four operators follows from an analysis of what operations
on possibility space are conceptually irreducible.

Pop is irreducible because commitment---the foreclosure of alternative
futures---cannot be expressed as a combination of documentation, dependency
recording, or observation. It is the primitive act of narrowing possibility.

Refuse is irreducible because documented inadmissibility is distinct from
commitment: it records that a path was considered and found inadmissible, which
is information that commitment cannot record. A commit of the alternative does not
document why the refused path was refused. Reasons are evidence that survives
collapse; without them, refusal is indistinguishable from non-consideration.

Bind is irreducible because dependency recording---coupling two trajectory
elements---cannot be expressed as commitment or refusal. A bind creates a
relationship between elements without consuming either; no combination of Pop
and Refuse produces this effect.

Collapse is irreducible because observation under a specific rule is distinct from
all three of the above. Pop, Refuse, and Bind all operate on the history and
option space; Collapse maps from history space to observable state space under
a specified projection. It is the meta-operation that makes histories visible.

\begin{theorem}[Completeness of the Four Operators]
\label{thm:completeness}
The four operators $\{P_x, R_{x,r}, B_{a,b}, C_{x,c}\}$ are sufficient to express:
variable assignment, conditional branching, function application, exception
handling, memory deallocation, and concurrent event recording.
\end{theorem}

\begin{proof}
We establish expressibility for each construct.

\textit{Variable assignment $x := v$}: Express as $P_v$ followed by $B_{x,v}$.
The commitment records the value choice; the bind records the name-value
dependency.

\textit{Conditional branch on predicate $\phi$}: Express as $P_{\phi^+}$ if $\phi$
holds, or $R_{\phi^+,\mathsf{ConstraintViolation}(\phi)}$ followed by
$P_{\phi^-}$ if not. The refusal records why the positive branch was not taken.

\textit{Function application $f(a)$}: Expressed as $\mathsf{LamIntro}(x)$ followed
by $\mathsf{Apply}(a)$---a deferred Pop. The closure freezes the option space at
the moment of abstraction; application is the moment of commitment.

\textit{Exception throwing}: Expressed as $R_{t, r}$ (documented inadmissibility
without system destruction). The type $\Rfsd(T, r)$ carries the reason as a
first-class certificate.

\textit{Memory deallocation}: Expressed as $R_{p, \mathsf{Deallocated}}$ where
$p$ is the pointer symbol. The refusal records that the memory region is no longer
admissible to access, making use-after-free a type-level error.

\textit{Concurrent event recording}: Expressed via Meld, the monoidal composition
of two histories. Two concurrent worlds $(H_1, \Omega_1)$ and $(H_2, \Omega_2)$
are combined as $(H_1 \mathbin{++} H_2, \Omega_1 \cup \Omega_2)$ under an
appropriate ordering convention.
\end{proof}

\section{The Conservation Law}

The four operators satisfy a conservation law that reflects the fundamental
structure of possibility.

\begin{theorem}[Conservation of Possibility]
\label{thm:conservation}
For any legal execution sequence
\[
W_0 \xrightarrow{e_1} W_1 \xrightarrow{e_2} \cdots \xrightarrow{e_n} W_n
\]
where each step is one of $P_x$, $R_{x,r}$, $B_{a,b}$, or $C_{x,c}$:
\begin{enumerate}
\item $|H_t|$ is strictly monotonically increasing: $|H_{t+1}| = |H_t| + 1$.
\item $|\Omega_t|$ is monotonically non-increasing.
\item Under pure Pop dynamics, $|H_t| + |\Omega_t| = |\Omega_0|$ for all $t$.
\end{enumerate}
\end{theorem}

\begin{proof}
(i) Every operator appends exactly one event to $H$, so $|H_{t+1}| = |H_t| + 1$.

(ii) $P_x$ removes $x$ from $\Omega$; the remaining operators do not modify
$\Omega$. Hence $|\Omega_{t+1}| \leq |\Omega_t|$.

(iii) Under pure Pop dynamics, every step removes exactly one element from $\Omega$
and appends one event to $H$. Starting from $|H_0| = 0$ and $|\Omega_0|$, after
$t$ steps $|H_t| = t$ and $|\Omega_t| = |\Omega_0| - t$, giving
$|H_t| + |\Omega_t| = |\Omega_0|$.
\end{proof}

The conservation law has the structure of a physical conservation principle:
possibility is not created or destroyed---it is either consumed (by Pop) or
documented as inadmissible (by Refuse). The operators Bind and Collapse record
relationships and observations without affecting the possibility budget.

\begin{definition}[Generalised Possibility Functional]
Define the event weight function $w : E \to \mathbb{N}$ by $w(\Pop(x)) = 1$,
$w(\Refuse(x,r)) = 0$, $w(\Bind(a,b)) = 0$, $w(\Collapse(x,c)) = 0$. For a
world $(H, \Omega)$, the \emph{generalised possibility functional} is
\[
\Pi(H, \Omega) = |\Omega| + \sum_{e \in H} w(e).
\]
\end{definition}

\begin{theorem}[Conservation of the Possibility Functional]
For any legal execution sequence, $\Pi(H_t, \Omega_t) = |\Omega_0|$ for all $t$.
\end{theorem}

\begin{proof}
At each step, exactly one of four operators is applied. Pop: $|\Omega|$ decreases
by 1; $w(\Pop) = 1$ is added. Net change: $(-1) + 1 = 0$. Refuse, Bind, Collapse:
$|\Omega|$ unchanged; $w = 0$ is added. Net change: $0$. In every case
$\Pi(H_{t+1}, \Omega_{t+1}) = \Pi(H_t, \Omega_t)$.
\end{proof}

\begin{corollary}[Irreversibility of Execution]
For any non-empty execution sequence, $|H_{t+1}| > |H_t|$. Therefore no legal
execution can return to a previous world state $(H_t, \Omega_t)$.
\end{corollary}

\begin{proof}
Every operator appends exactly one event to $H$, so $|H_{t+1}| = |H_t| + 1 > |H_t|$.
Since $H$ grows strictly, $(H_{t+1}, \Omega_{t+1}) \neq (H_t, \Omega_t)$.
\end{proof}

This irreversibility is not a limitation but a feature. It is the formal
correlate of the phenomenological observation that histories are primary: a world
that could return to a prior state would be a world in which the history of
the intervening steps was undone, which is not computation but its negation.

\chapter{The Algebra of Histories}

\section{Histories as the Primary Objects}

Before introducing types, evaluators, or compilers, we must characterize the
mathematical structure of the objects that are primary in the Spherepop ontology.
Histories are not merely lists of events. They form a category, and the category
structure constrains what operations on histories are meaningful.

\begin{definition}[History]
A \emph{history} over event alphabet $E$ is a finite sequence
$H = [e_1, e_2, \ldots, e_n]$ with each $e_i \in E$. The empty history is
$\varepsilon = []$.
\end{definition}

\begin{definition}[History Concatenation]
For histories $H_1 = [e_1, \ldots, e_m]$ and $H_2 = [f_1, \ldots, f_n]$,
\[
H_1 \mathbin{++} H_2 = [e_1, \ldots, e_m, f_1, \ldots, f_n].
\]
\end{definition}

\begin{proposition}[History Monoid]
The triple $(\calH, \mathbin{++}, \varepsilon)$ is the free monoid over $E$.
\end{proposition}

\begin{proof}
Associativity of $\mathbin{++}$ follows from list concatenation associativity.
$\varepsilon$ is the left and right identity. Freeness: every history $H$ is
uniquely determined by its event sequence, and no non-trivial relation among
generators holds.
\end{proof}

\section{The History Category}

\begin{definition}[History Category $\Hist$]
Define the category $\Hist$ as follows:
\begin{itemize}
\item \textit{Objects}: option spaces $\Omega \subseteq \Omega_0$.
\item \textit{Morphisms}: $\mathrm{Hom}(\Omega, \Omega')$ is the set of histories
$H$ such that starting from $\Omega$, executing $H$ yields $\Omega'$.
\item \textit{Composition}: $H_2 \circ H_1 = H_1 \mathbin{++} H_2$ (execute
$H_1$ first, then $H_2$).
\item \textit{Identity}: $\mathrm{id}_\Omega = \varepsilon$.
\end{itemize}
\end{definition}

\begin{proposition}[$\Hist$ is a Well-Defined Category]
$\Hist$ satisfies all category axioms.
\end{proposition}

\begin{proof}
Associativity of composition follows from associativity of $\mathbin{++}$.
$\varepsilon$ is the identity morphism since $H \mathbin{++} \varepsilon = H =
\varepsilon \mathbin{++} H$.
\end{proof}

\begin{proposition}[$\Hist$ is a Free Strict Monoidal Category]
$(\Hist, \mathbin{++}, \varepsilon, \otimes)$ is a free strict monoidal category
where the tensor product $H_1 \otimes H_2$ is the Meld operation (parallel
composition of histories) and the monoidal unit is the empty option space with
empty history.
\end{proposition}

\begin{remark}
The free monoid structure forces \texttt{append} (extend by one generator) and
\texttt{meld} (monoidal composition of two histories) as the only legitimate
history operations. The absence of \texttt{remove} and \texttt{undo} is the
mechanization of freeness: there are no relations among generators.
\begin{lstlisting}[style=rust, caption={History as free monoid in Rust}]
pub fn append(&mut self, event: Event) {
    self.events.push(event);  // one generator
}
pub fn meld(&mut self, other: &History) {
    self.events.extend(other.events.iter().cloned());
}
\end{lstlisting}
History equality $H_1 = H_2$ is the primary criterion of program identity.
\end{remark}

\section{Unique Factorisation}

\begin{theorem}[Unique History Factorisation]
Every non-empty history $H = [e_1, e_2, \ldots, e_n]$ admits a unique
factorisation into primitive generators:
\[
H = [e_1] \mathbin{++} [e_2] \mathbin{++} \cdots \mathbin{++} [e_n].
\]
No other factorisation into generators exists.
\end{theorem}

\begin{proof}
Since $\calH$ is the free monoid over $E$, every element has a unique normal form
as a sequence of generators. The generators are singletons $[e]$; freeness of the
monoid guarantees there are no non-trivial relations among generators.
\end{proof}

\begin{remark}
Unique factorisation is what makes history equivalence a decidable correctness
criterion for compilation. If two execution paths produce the same history, they
must have produced the same sequence of primitive events in the same order.
\end{remark}

\section{Admissibility as a Set-Valued Function}

\begin{definition}[Admissible Futures]
For a history $H$, the set of admissible future symbols is
\[
\calA(H) = \Omega_0 \setminus \{x \mid \exists r.\, \Refuse(x, r) \in H\}.
\]
\end{definition}

\begin{definition}[Admissibility Contraction]
Let $H' = H \mathbin{++} [\Refuse(x, r)]$. Then $\calA(H') \subseteq \calA(H)$,
and specifically $x \notin \calA(H')$ while $x$ may have been in $\calA(H)$.
\end{definition}

\begin{remark}
The asymmetry between Pop and Refuse is critical. Pop contracts the structural
option space $\Omega$. Refuse contracts the admissibility set $\calA(H)$ without
touching $\Omega$. Two execution paths may have the same $\Omega$ while having
different $\calA(H)$---they are structurally equivalent but semantically
distinguished by their refusal histories. This distinction is invisible to any
state-centric computational model.
\end{remark}

%% ============================================================
%% PART III: OPERATIONAL SEMANTICS AND COLLAPSE
%% ============================================================
\part{Operational Semantics and Collapse}

\chapter{Small-Step Semantics}

\section{The Transition System}

We write $(H, \Omega) \xrightarrow{\alpha} (H', \Omega')$ for a single execution
step labelled $\alpha$.

\[
\frac{x \in \Omega}{(H, \Omega) \xrightarrow{\mathsf{pop}(x)} (H \mathbin{++} [\Pop(x)],\; \Omega \setminus \{x\})} \quad [\mathrm{POP}]
\]

\[
\frac{}{(H, \Omega) \xrightarrow{\mathsf{refuse}(x,r)} (H \mathbin{++} [\Refuse(x, r)],\; \Omega)} \quad [\mathrm{REFUSE}]
\]

Note: $[\mathrm{REFUSE}]$ has no premise---any symbol may be refused. The effect
on the option space is nil; the effect on admissibility is non-nil.

\[
\frac{}{(H, \Omega) \xrightarrow{\mathsf{bind}(a,b)} (H \mathbin{++} [\Bind(a, b)],\; \Omega)} \quad [\mathrm{BIND}]
\]

\[
\frac{\Gamma \vdash t : \Adm(T)}{(H, \Omega) \xrightarrow{\mathsf{collapse}(x,c)} (H \mathbin{++} [\Collapse(x, c)],\; \Omega)} \quad [\mathrm{COLLAPSE}]
\]

The premise $\Gamma \vdash t : \Adm(T)$ is the type-theoretic admissibility guard.
Collapse requires an admissibility certificate. Without it, collapse is a hard
type error.

\section{Big-Step Semantics}

The big-step relation $\langle t, W \rangle \Downarrow \langle v, W' \rangle$
relates terms and worlds to values and resulting worlds.

\begin{definition}[Evaluation Result]
An \emph{evaluation result} is a triple $(v, H, a)$ where $v$ is a value,
$H$ is the accumulated history, and $a \in \{\top, \bot\}$ is the admissibility
flag. The admissibility flag is $\top$ if every operation in $H$ was admissible
and $\bot$ if any refusal was encountered.
\end{definition}

The key invariant is that evaluation never discards history. Every step appends
to $H$; no step modifies or removes prior events. The world grows monotonically.

\begin{lstlisting}[style=rust, caption={Evaluation result preserving all history}]
pub struct EvalResult {
    pub value: Value,
    pub world: World,    // world.history grows monotonically
    pub admissible: bool,
}
\end{lstlisting}

\section{Multi-Step Reduction}

\begin{definition}[Multi-Step Reduction]
The multi-step reduction $\to^*$ is the reflexive transitive closure of $\to$.
A term $t$ is \emph{normal} if there is no $t'$ with $t \to t'$.
\end{definition}

\begin{theorem}[History Monotonicity]
If $(H_0, \Omega_0) \to^* (H_n, \Omega_n)$ then $H_0$ is a prefix of $H_n$:
there exists $H'$ such that $H_n = H_0 \mathbin{++} H'$.
\end{theorem}

\begin{proof}
By induction on the length of the reduction sequence. Each step appends exactly
one event to the history. The result follows by transitivity.
\end{proof}

\begin{theorem}[Option Space Monotonicity]
If $(H_0, \Omega_0) \to^* (H_n, \Omega_n)$ then $\Omega_n \subseteq \Omega_0$.
\end{theorem}

\begin{proof}
Only Pop removes elements from $\Omega$; all other operators leave it unchanged.
By induction, $\Omega$ can only shrink.
\end{proof}

\chapter{Collapse as Quotient}

\section{The Insufficiency of Unparameterized Collapse}

The original Spherepop design represented collapse as simply \texttt{Collapse(Symbol)}.
This was insufficient, and the insufficiency is philosophically illuminating.

Observable state depends on how you look. The same history, examined by different
observers with different observational frameworks, yields different observable
states. A collapse without a rule is a collapse into an unspecified notion of
visibility. The rule is not an implementation detail; it is what makes the
observation meaningful.

\begin{definition}[Collapse Rule]
A \emph{collapse rule} is a function $c : \calH \to O_c$ from the history monoid
to an observational space $O_c$.
\end{definition}

\begin{definition}[Observational Equivalence]
Two histories $H_1, H_2 \in \calH$ are \emph{observationally equivalent under
rule $c$}, written $H_1 \sim_c H_2$, if and only if $c(H_1) = c(H_2)$.
\end{definition}

\begin{definition}[Observable State as Quotient]
The \emph{observable state space} under rule $c$ is the quotient:
\[
O_c = \calH / {\sim_c}.
\]
An observable state is an equivalence class of histories.
\end{definition}

Observable state is not a state. Observable state is a quotient. Two histories
that are globally distinct may be observationally indistinguishable under a
specific collapse rule. This is a feature: the collapse rule specifies which
distinctions are relevant to a given question.

\section{The Four Canonical Rules}

\begin{definition}[Identity Rule]
$c_I : \calH \to \calH$ is the identity: $c_I(H) = H$. Under the identity rule,
every event is visible and $O_I \cong \calH$.
\end{definition}

\begin{definition}[LastWrite Rule]
$c_{LW}$ maps $H$ to the most recent pop for each symbol, with subsequent refusals
retroactively removing symbols:
\[
c_{LW}(H) = \{x \mapsto 1 \mid \Pop(x) \in H \text{ and } \nexists r.\, \Refuse(x,r) \in H \text{ after the last } \Pop(x)\}.
\]
\end{definition}

\begin{definition}[Accumulate Rule]
$c_{\mathrm{acc}}(H)$ maps each symbol to the count of its Pop events:
\[
c_{\mathrm{acc}}(H)(x) = |\{i \mid e_i = \Pop(x)\}|.
\]
\end{definition}

\begin{definition}[Projection Rule]
For a label prefix $L$, the projection rule $\pi_L$ restricts to events whose
symbol names share prefix $L$:
\[
\pi_L(H) = [e_i \in H \mid \mathrm{sym}(e_i) \text{ starts with } L].
\]
\end{definition}

\begin{proposition}[Coarsening Order]
The collapse rules form a partial order by coarsening: $c_1 \leq c_2$ iff
${\sim_{c_2}} \subseteq {\sim_{c_1}}$ (finer equivalence = more informative
observation). Under this order: $c_I$ is the finest; $\pi_L$ is coarser than
$c_I$; $c_{LW}$ and $c_{\mathrm{acc}}$ are incomparable in general.
\end{proposition}

\section{The Collapse Functor}

\begin{definition}[Observable-State Category $\Obs_c$]
For a collapse rule $c$, define the category $\Obs_c$:
\begin{itemize}
\item \textit{Objects}: observable states $o \in O_c$.
\item \textit{Morphisms}: $\mathrm{Hom}(o_1, o_2)$ is the set of observable
transitions transforming $o_1$ into $o_2$ under rule $c$.
\item \textit{Composition}: sequential observable transition.
\end{itemize}
\end{definition}

\begin{definition}[Collapse Functor]
For collapse rule $c$, define $F_c : \Hist \to \Obs_c$ by
$F_c(\Omega) = c(\varepsilon_\Omega)$ on objects and $F_c(H) = c(H)$ on morphisms.
\end{definition}

\begin{theorem}[Functoriality of Collapse]
$F_c$ is a well-defined functor when $c$ respects composition:
\[
c(H_2 \mathbin{++} H_1) = c(H_2) \circ c(H_1).
\]
\end{theorem}

\begin{proof}
Identity preservation: $F_c(\varepsilon) = c(\varepsilon) = \mathrm{id}_{O_c}$.
Composition preservation: $F_c(H_2 \circ H_1) = F_c(H_1 \mathbin{++} H_2) =
c(H_1 \mathbin{++} H_2) = c(H_2) \circ c(H_1) = F_c(H_2) \circ F_c(H_1)$ when
$c$ respects composition.
\end{proof}

\begin{remark}
The functoriality of collapse elevates it from an implementation detail to a
category-theoretic object. Collapse is a structure-preserving map between
categories. The collapse rule specifies which structures are preserved, and
different rules preserve different structures. $c_{LW}$ is not generally functorial
because the last-write of a concatenated history is not the composition of the
last-writes of its parts---a later segment can retroactively affect the observable
state of an earlier one.
\end{remark}

\section{The Universal Property of Collapse}

\begin{theorem}[Universal Property of Collapse]
Let $q_c : \calH \to O_c$ be the quotient map for rule $c$. If $f : \calH \to X$
is any function satisfying
\[
H_1 \sim_c H_2 \implies f(H_1) = f(H_2),
\]
then there exists a unique $\bar{f} : O_c \to X$ such that $f = \bar{f} \circ q_c$.
\end{theorem}

\begin{proof}
By the universal property of quotient sets. Define $\bar{f}([H]_c) = f(H)$.
This is well-defined because $f$ is constant on equivalence classes. Uniqueness:
any $\bar{f}$ satisfying $f = \bar{f} \circ q_c$ must map $[H]_c$ to $f(H)$.
\end{proof}

\begin{remark}
The universal property says that observable state is not one representation among
many. It is the universal representation among all representations that identify
the same histories as $c$ does. Every coarser observation necessarily passes
through $O_c$.
\end{remark}

%% ============================================================
%% PART IV: TYPE THEORY
%% ============================================================
\part{Type Theory}

\chapter{Types as Admissibility Certificates}

\section{The Inadequacy of Value-Set Types}

The standard interpretation of types---a type is a set of admissible values---is
adequate for state-centric languages. It is insufficient for Spherepop.

Consider what a type must certify in a history-primary setting. A value $v$ of
type $T$ is not merely a member of a set. It is the result of a path through
history that reached $v$ without violating any constraint represented by $T$.
The type is not a description of $v$ at a moment; it is a certificate about the
history that produced $v$.

\section{The Type Language}

\begin{definition}[Spherepop Types]
The type language of Spherepop is defined inductively:
\begin{align*}
T ::=\; & \mathsf{Unit} \mid \mathsf{Never} \mid X \\
\mid\; & \Pi(x : T_1).T_2 \\
\mid\; & \Adm(T) \\
\mid\; & \Rfsd(T, r) \\
\mid\; & \Col(T, c) \\
\mid\; & \Proc(T_1 \rightsquigarrow T_2)
\end{align*}
The simple function type $T_1 \to T_2$ is notation for $\Pi(\_ : T_1).T_2$.
\end{definition}

The distinctive types carry certificates, not just values:
\begin{itemize}
\item $\Adm(T)$: the path that produced this term was admissible.
\item $\Rfsd(T, r)$: the term was encountered but found inadmissible for reason
$r$. The reason is part of the type---two terms refused for different reasons
have genuinely different types.
\item $\Col(T, c)$: an admissible term was observed under collapse rule $c$ and
the observation is sealed. The rule $c$ is part of the type.
\end{itemize}

\section{The Type Rules}

We write $\Gamma \vdash t : T$ for the standard typing judgement.

\[
\frac{\Gamma \vdash t : T}{\Gamma \vdash \mathsf{pop}(t) : \Adm(T)} \quad [\mathrm{T\text{-}POP}]
\]

Pop transforms any typed term into an admissible term. It is the mechanism by
which structural commitment produces a reachability certificate.

\[
\frac{\Gamma \vdash t : T \quad r \in \mathsf{RefusalReason}}{\Gamma \vdash \mathsf{refuse}(t, r) : \Rfsd(T, r)} \quad [\mathrm{T\text{-}REFUSE}]
\]

Refusal is typed: the type carries the reason. A term refused for
$\mathsf{ConstraintViolation}(c)$ has a type formally distinct from one refused
for $\mathsf{NotAvailable}$.

\[
\frac{\Gamma \vdash t : \Adm(T) \quad c \in \mathsf{CollapseRule}}{\Gamma \vdash \mathsf{collapse}(t, c) : \Col(T, c)} \quad [\mathrm{T\text{-}COLLAPSE}]
\]

Collapse requires an admissibility certificate. A term whose type is not
$\Adm(\cdot)$ cannot be collapsed---this is a hard type error, not a runtime
exception.

\[
\frac{\Gamma \vdash a : T_1 \quad \Gamma \vdash b : T_2}{\Gamma \vdash \mathsf{bind}(a, b) : \Proc(T_1 \rightsquigarrow T_2)} \quad [\mathrm{T\text{-}BIND}]
\]

\[
\frac{\Gamma, x : T_1 \vdash b : T_2}{\Gamma \vdash \lambda x : T_1.\, b : \Pi(x : T_1).T_2} \quad [\mathrm{T\text{-}LAM}]
\]

\[
\frac{\Gamma \vdash f : \Pi(x : T_1).T_2 \quad \Gamma \vdash a : T_1}{\Gamma \vdash f\, a : T_2[x := a]} \quad [\mathrm{T\text{-}APP}]
\]

\section{Open Worlds and Closed Worlds}

During implementation of the Spherepop type checker, a conflict emerged between
how the interpreter and type checker treated free variables. The apparent bug
conceals a genuine theoretical discovery: which behavior is correct depends on
which assumption about the boundary of the world is in force.

\begin{definition}[Closed-World Context]
A \emph{closed-world context} $\Gamma_c$ satisfies:
\[
x \notin \Gamma_c \implies \Gamma_c \vdash x : \bot.
\]
The absence of a declaration is the absence of the object.
\[
\frac{x \notin \Gamma \quad \Gamma \text{ is closed-world}}{\Gamma \not\vdash x : T \text{ for any } T} \quad [\mathrm{T\text{-}VAR\text{-}CLOSED}]
\]
\end{definition}

\begin{definition}[Open-World Context]
An \emph{open-world context} $\Gamma_o$ satisfies:
\[
x \notin \Gamma_c \implies \Gamma_o \vdash x : \mathsf{Unit}.
\]
The absence of a declaration is undecided possibility.
\[
\frac{x \notin \Gamma \quad \Gamma \text{ is open-world}}{\Gamma \vdash x : \mathsf{Unit}} \quad [\mathrm{T\text{-}VAR\text{-}OPEN}]
\]
\end{definition}

\begin{proposition}[Monotonicity of Provability]
$\Gamma_c \subseteq \Gamma_o \implies \mathrm{Provable}(\Gamma_c) \subseteq
\mathrm{Provable}(\Gamma_o)$.
\end{proposition}

\begin{remark}
The open/closed world distinction is not an engineering convenience. It is the
mechanization of an epistemological commitment: a closed-world type checker assumes
complete knowledge of the reachable; an open-world type checker assumes partial
knowledge. A language whose fundamental ontology is about possibility space cannot
leave this distinction implicit. The $\mathsf{TypeMode}$ is preserved through
\texttt{extend}: a closed-world context cannot silently become open-world
mid-derivation.
\end{remark}

\section{Principal Theorems}

\begin{theorem}[Collapse Soundness]
If $\Gamma \vdash \mathsf{collapse}(t, c) : T$ then $\Gamma \vdash t : \Adm(T')$
for some $T'$, and $T = \Col(T', c)$.
\end{theorem}

\begin{proof}
By inversion on $[\mathrm{T\text{-}COLLAPSE}]$: the only rule that derives a
Collapse judgement requires the premise $\Gamma \vdash t : \Adm(T)$, so the
inner term must have an admissible type.
\end{proof}

\begin{theorem}[Type Preservation]
If $\Gamma \vdash t : T$ and $t \to t'$, then $\Gamma \vdash t' : T$.
\end{theorem}

\begin{proof}
By structural induction on the reduction relation $\to$. The key cases are
$[\mathrm{T\text{-}POP}]$, $[\mathrm{T\text{-}REFUSE}]$, and
$[\mathrm{T\text{-}COLLAPSE}]$, each mapping a reduction step to an admissibility
transformation closed under the respective rule. Beta-reduction requires a
substitution lemma: $\Gamma, x : T_1 \vdash b : T_2$ and $\Gamma \vdash a : T_1$
implies $\Gamma \vdash b[x := a] : T_2[x := a]$.
\end{proof}

\begin{theorem}[Progress]
If $\vdash t : T$ (well-typed in the empty context) then either $t$ is a value,
or there exists $t'$ such that $t \to t'$.
\end{theorem}

\begin{proof}
By structural induction on $t$. Variable terms are values. Lambda terms are values
(closures). $\mathsf{pop}$, $\mathsf{refuse}$, $\mathsf{bind}$, and
$\mathsf{collapse}$ of a well-typed inner term either reduce (if the inner term
reduces) or produce a value (if the inner term is a value).
\end{proof}

\chapter{Dependent Types and Admissibility Lattices}

\section{Dependent Process Types}

The current $\Proc(T_1 \rightsquigarrow T_2)$ is first-order. A dependent process
type allows the output type to depend on the specific value consumed.

\begin{definition}[Dependent Process Type]
A \emph{dependent process type}
\[
\Pi(x : T_1).\, \Proc(x, f(x))
\]
allows the type of the output process to depend on the value $x$ of the input.
This is the natural extension of dependent type theory into the process calculus
setting.
\end{definition}

\begin{example}
A function that reads a file and returns a proof that the file was read can be
typed as:
\[
\Pi(\mathit{path} : \mathsf{FilePath}).\, \Proc(\mathit{path}, \Adm(\mathsf{FileContents}(\mathit{path}))).
\]
The $\Adm$ in the output type certifies that the reading path was admissible---no
access violation occurred, the file existed, and the permissions were valid.
\end{example}

\section{Admissibility Lattices}

The current admissibility structure is Boolean: a term is admissible or not. A
more refined theory recognizes that admissibility is graded.

\begin{definition}[Admissibility Lattice]
An \emph{admissibility lattice} $(\mathcal{L}, \leq, \top, \bot, \wedge, \vee)$
is a bounded lattice where:
\begin{itemize}
\item $\top$ is full admissibility: the trajectory satisfies all constraints.
\item $\bot$ is total inadmissibility: the trajectory violates all constraints.
\item $\ell_1 \wedge \ell_2$ is the admissibility under the conjunction of
constraints $c_1$ and $c_2$.
\item $\ell_1 \vee \ell_2$ is the admissibility under the disjunction.
\end{itemize}
\end{definition}

\begin{definition}[Graded Admissibility Type]
$\Adm_\ell(T)$ denotes admissibility at level $\ell \in \mathcal{L}$. The
standard $\Adm(T)$ is $\Adm_\top(T)$.
\end{definition}

\begin{proposition}[Admissibility Lattice Typing]
The typing rule for graded admissibility is:
\[
\frac{\Gamma \vdash t : T \quad \ell \in \mathcal{L}}{\Gamma \vdash \mathsf{pop}_\ell(t) : \Adm_\ell(T)}
\]
Collapse requires admissibility at level at least $\ell_0$ for some threshold
$\ell_0$ determined by the collapse rule $c$.
\end{proposition}

\begin{remark}
Admissibility lattices connect the Spherepop type theory to the admissibility
geometry program of the RSVP and CLIO frameworks. The xylomorphic criterion
$\lambda < 1$ from the admissibility field is a threshold condition in a
continuous admissibility lattice, where admissibility is measured by a scalar
$\lambda$ and collapse is only valid when $\lambda$ falls below the threshold.
\end{remark}

\section{Proof-Carrying Refusal}

The $\mathsf{RefusalReason}$ type is currently a vocabulary for inadmissibility.
It is not yet a proof system. A fully realized Spherepop type system would require
$\mathsf{RefusalReason}$ to be a proof term: evidence that can be checked, not
merely a label.

\begin{definition}[Proof-Carrying Refusal]
A \emph{proof-carrying refusal} is a refusal of the form
$\mathsf{refuse}(t, \pi)$ where $\pi$ is a proof term of type
$\mathsf{Inadmissible}(t, \Omega)$---a formal certificate that $t$ is inadmissible
given the current admissibility set $\calA(H)$.
\end{definition}

\begin{definition}[Inadmissibility Proof System]
The inadmissibility propositions and their proofs are:
\begin{align*}
\mathsf{AlreadyRefused}(x) &: x \in \{y \mid \exists r.\, \Refuse(y,r) \in H\} \\
\mathsf{ConstraintViolation}(c, x) &: c(x) > \mathsf{threshold}(c) \\
\mathsf{NotAvailable}(x) &: x \notin \Omega \\
\mathsf{DependencyFailed}(x, y) &: \Bind(x,y) \in H \wedge \mathsf{AlreadyRefused}(y)
\end{align*}
Each proposition has a type-theoretic proof term, and $\mathsf{refuse}(t, \pi)$
is only well-typed when $\pi$ proves one of these propositions for $t$.
\end{definition}

\begin{theorem}[Refusal Completeness]
Every inadmissible trajectory has a proof-carrying refusal at some point.
That is, if $\calA(H) \not\ni x$ for some $x$ that is later attempted, then
there exists a proof term $\pi$ of $\mathsf{Inadmissible}(x, \calA(H))$.
\end{theorem}

\begin{proof}
$x \notin \calA(H)$ implies $\exists r.\, \Refuse(x, r) \in H$. The proof term
$\pi = \mathsf{AlreadyRefused}(x)$ witnesses that $x$ is in the refused set,
which is defined as the complement of $\calA(H)$.
\end{proof}

%% ============================================================
%% PART V: NOTES AS SPHEREPOP OBJECTS
%% ============================================================
\part{Notes as Spherepop Objects}

\chapter{The Seven Note Classes as Operator Realizations}

The four Spherepop operators---Pop, Refuse, Bind, Collapse---are abstract
operations on possibility space. Notes are their physical realization in cognition,
computation, language, and civilization. This chapter develops the correspondence
systematically.

\section{The Classification Principle}

\begin{definition}[Dominant Operator]
The \emph{dominant operator} of a note class is the Spherepop operator that
accounts for the primary function of notes in that class. Most notes mix operators,
but each class has a dominant one.
\end{definition}

\begin{theorem}[Notes as Spherepop Programs]
Every note is a Spherepop program. Every Spherepop program is a note. The
correspondence is given by the following bijection between note classes and
operator combinations:
\begin{center}
\begin{tabular}{lll}
\toprule
Note Class & Dominant Operator & Secondary Operators \\
\midrule
Capture & Refuse & --- \\
Prospective & Bind & Pop (deferred) \\
Repair & Refuse $+$ Pop & --- \\
Coordination & Bind & Refuse \\
Refusal & Refuse & --- (pure form) \\
Generative & Pop & Bind \\
Collapse & Collapse & --- \\
\bottomrule
\end{tabular}
\end{center}
\end{theorem}

\chapter{Capture Notes}

\begin{definition}[Capture Note]
A \emph{capture note} is a note whose primary function is to preserve a
continuation that exists at the time of note creation but is expected to become
inaccessible. In Spherepop terms, a capture note performs a Refuse operation
with respect to the projection $\sigma : \calH \to S$: it preserves history
fibers that the projection would collapse.
\end{definition}

A photograph refuses the collapse of a perceptual trajectory. An audio recording
refuses the collapse of a sonic trajectory. A laboratory notebook refuses the
collapse of a procedural and cognitive trajectory.

\begin{proposition}[Capture Notes and Refuse]
A capture note $N$ performs a Refuse operation with respect to the projection
$\sigma : \calH \to S$ if and only if there exists a history fiber
$h \in \sigma^{-1}(S)$ such that $P_A(c_h, t \mid N) > P_A(c_h, t)$, where
$c_h$ is the continuation associated with history $h$.
\end{proposition}

The quality criterion for a capture note is not fidelity to the original state
but adequacy as a reconstruction substrate. A photograph that captures the correct
person but not the context may be a poor capture note even at high resolution, if
the lost context is what matters for future reconstruction. The adequacy of a
capture note depends on what future continuations are anticipated.

\chapter{Prospective Notes}

\begin{definition}[Prospective Note]
A \emph{prospective note} is a note whose primary function is to preserve access
to a continuation that does not yet exist at note creation time but is anticipated.
In Spherepop terms, a prospective note performs a Bind operation between present
agent state $A_t$ and future continuation $c_{t'}$.
\end{definition}

A calendar entry binds future action to present scheduling. A blueprint binds
future construction to present design. A roadmap binds future development to
present strategic commitment. Source code binds future computation to present
specification---and is unique in that the binding is mechanically executable.

\begin{proposition}[Prospective Notes and Temporal Binding]
A prospective note $N$ performs a Bind operation if and only if
\[
P_A(c_{t'}, t' \mid N, A_t) > P_A(c_{t'}, t' \mid A_t),
\]
where $t' > t$ and the conditioning on $A_t$ indicates the agent's present state
is available in both cases.
\end{proposition}

\begin{proposition}[Collective Reachability from Coordination]
Let $A_1, \ldots, A_n$ be agents with independent continuation spaces
$\calC_1, \ldots, \calC_n$. A coordination note $N$ increases collective
reachability if
\[
|\calC(A_1, \ldots, A_n \mid N)| > |\calC_1| + \cdots + |\calC_n|.
\]
\end{proposition}

The proposition captures the supraadditive character of coordination: bound agents
can traverse paths that no individual could reach alone.

\chapter{Repair, Refusal, and Generative Notes}

\section{Repair Notes}

\begin{definition}[Repair Note]
A \emph{repair note} is a note whose primary function is to reduce the cost of
reconstructing a continuation after it has been or is anticipated to be
interrupted. It combines Refuse (preserving the history of the trajectory against
collapse) and Pop (enabling re-entry at a specific point after interruption).
\end{definition}

\begin{theorem}[Notes as Pre-emptive Repair]
Every note $N$ with positive repair value $R_v(N, c) = C_r(c) - C_r(c \mid N) > 0$
constitutes a pre-emptive repair: it reduces expected reconstruction cost prior
to any failure occurring.
\end{theorem}

\begin{proof}
$R_v(N, c) > 0$ implies $C_r(c \mid N) < C_r(c)$. Since the note exists before
any failure occurs, the cost reduction is in force from the moment of note
creation. The note therefore functions as a repair mechanism that precedes the
event it repairs for.
\end{proof}

A troubleshooting guide is a repair note: a stored Pop operator for a class of
failure states. Each entry says: if you arrive at this failure state, the
continuation from here to a functional state runs through these steps.

\section{Refusal Notes}

\begin{definition}[Refusal Note]
A \emph{refusal note} is a note whose primary function is to prevent the collapse
of an existing distinction, independently of whether the preserved distinction is
expected to be used. It is the closest realization of the pure Refuse operation.
\end{definition}

Mathematical proofs, cryptographic checksums, archival copies, legal records, and
formal specifications belong to this class.

\begin{proposition}[Proofs as Refusal Notes]
A mathematical proof $P$ of theorem $T$ from premises $\Gamma$ satisfies
\[
P_A(c_{\Gamma \vdash T}, t \mid P) > P_A(c_{\Gamma \vdash T}, t)
\]
for any agent $A$ with sufficient background knowledge, at any future time $t$.
\end{proposition}

The proof refuses the collapse of an inferential trajectory into mere assertion.
Its value is not actualized use but preserved potential: it maintains the
reachability of the derivation whether or not anyone consults it.

\section{Generative Notes}

\begin{definition}[Generative Note]
A \emph{generative note} is a note that creates access to continuations that were
not reachable before the note's creation. It is the primary realization of the Pop
operation: it opens new trajectories in possibility space.
\end{definition}

\begin{proposition}[Generative Notes and Novel Reachability]
A generative note $N$ satisfies
\[
\exists\, c \in \calC : P_A(c, t \mid N) > 0 \text{ and } P_A(c, t) = 0.
\]
That is, $N$ creates access to continuations that were previously unreachable.
\end{proposition}

Scientific theories, maps, programming languages, mathematical formalisms, and
ontologies are generative notes. A programming language is a generative note that
is mechanically executable: it defines not just a space of possible programs but
a machine for traversing that space.

\chapter{Collapse Notes and the Ephemerality Inversion}

\section{The Hidden Class}

\begin{definition}[Collapse Note]
A \emph{collapse note} is an artifact that performs a deliberate projection
$\sigma_N : \calH \to S_N$, compressing a history or artifact into a more compact
representation at the cost of detail, in order to make the representation
navigable. It is the realization of the Collapse operator.
\end{definition}

File names, titles, abstracts, indexes, summaries, taxonomies, classifications,
labels, and keywords are all collapse notes. The paradox of collapse notes is that
they appear to be the opposite of notes---they destroy information---yet they are
indispensable to the function of every other note class.

\begin{theorem}[Necessity of Collapse Notes]
For any sufficiently large collection of notes $\calN$, there exists a threshold
$|\calN| > K$ above which the reachability of individual notes in $\calN$ decreases
in the absence of collapse notes over $\calN$.
\end{theorem}

\begin{proof}
Search cost through an unstructured collection grows at least linearly with
collection size. For $|\calN| > K$, exhaustive search becomes too expensive for
practical use. A collapse note---index, catalogue, taxonomy---reduces search cost
by imposing structure, making individual notes reachable at sublinear cost.
Without such structure, notes become inaccessible through navigational failure,
not destruction.
\end{proof}

\begin{remark}
A word is a collapse note: it compresses a complex of experience, association,
inference, and usage into a single retrievable token. A noun is a collapse note
about a process. The passive theory would not recognize these as notes at all.
The Spherepop theory recognizes them immediately as the most pervasive and powerful
class of collapse notes civilization has developed.
\end{remark}

\section{The Ephemerality Inversion}

\begin{definition}[Continuation Completion]
A continuation $c$ is \emph{complete} with respect to note $N$ at time $t$ if
the trajectory encoded by $N$ has been successfully traversed.
\end{definition}

\begin{theorem}[Ephemerality Inversion]
Let $N$ be a note encoding a single continuation $c$. If $c$ has been completed,
then the destruction of $N$ does not reduce the reachability of any remaining
continuation. Therefore the ephemerality of $N$ after completion of $c$ implies
neither loss nor failure.
\end{theorem}

\begin{proof}
By assumption, $c$ is the sole continuation for which $N$ increases reachability.
After completion of $c$, $V_C(N) = 0$. Since continuation volume is zero, the
destruction of $N$ reduces reachability by zero.
\end{proof}

The grocery list discarded at the supermarket exit has not failed. It has succeeded
completely. The theorem generates an inversion of the passive theory's ranking:

\begin{center}
\begin{tabular}{lll}
\toprule
Fate & Passive Theory & Continuation Theory \\
\midrule
Executed and discarded & Low value & Optimal \\
Preserved and re-entered & Medium value & Valuable \\
Preserved, never re-entered & High value & Unrealized potential \\
\bottomrule
\end{tabular}
\end{center}

The celebrated survivors of antiquity---incomplete manuscripts, undelivered
letters, abandoned plans---survive precisely because their continuations were never
completed. The completed notes, the sent letters, the finished buildings, left no
trace because they succeeded.

%% ============================================================
%% PART VI: COMPILATION, GC, AND ERROR CORRECTION
%% ============================================================
\part{Compilation, Garbage Collection, and Error Correction}

\chapter{The Event IR and Compiler Correctness}

\section{The Failure of Traditional Compiler Correctness}

The standard notion of compiler correctness is: a compiler is correct if the
compiled program produces the same output as the interpreted program for every
input. This criterion is adequate for state-centric languages. It is insufficient
for Spherepop.

Two programs that produce the same final value may have constructed radically
different histories to get there, and those histories are part of what the program
is.

\begin{definition}[History Equivalence as Correctness]
Let $I : P \to \calH$ be the interpreter function mapping programs to histories,
and $V : \mathsf{IR} \to \calH$ the VM mapping event-IR blocks to histories,
and $C : P \to \mathsf{IR}$ the compiler. A compiler $C$ is correct if for all
programs $p$ and initial worlds $W$:
\[
I(p, W).\mathsf{history} = V(C(p), W).\mathsf{history}.
\]
Output equivalence (same final value) is a strictly weaker criterion that is
insufficient for Spherepop.
\end{definition}

\section{Event IR}

Conventional intermediate representations encode instructions that mutate machine
state. Spherepop IR encodes \emph{events}---records of occurrences that persist
as part of the history the VM is constructing.

\begin{definition}[Event IR]
The Spherepop intermediate representation is a sequence of $\mathsf{IrEvent}$
terms:
\[
\mathsf{IR} ::= \Pop \mid \Refuse(r) \mid \Collapse(c) \mid \Bind \mid \mathsf{Load}(x) \mid \mathsf{Apply} \mid \cdots
\]
Each event carries its reason or rule as a first-class argument, preserving the
full certificate structure of the source term.
\end{definition}

\begin{lstlisting}[style=rust, caption={Event IR in Rust}]
pub enum IrEvent {
    Pop,
    Refuse { reason: RefusalReason },
    Collapse { rule: CollapseRule },
    Bind,
    Load(Symbol),
    Apply,
    LamIntro(Symbol, Type),
}
\end{lstlisting}

\begin{proposition}[IR Faithfulness]
The compilation function $C : \mathsf{Term} \to \mathsf{IR}$ is faithful: for
every primitive term $t$, $C(t)$ contains exactly the event variants that $I(t)$
would append to the history, with identical arguments.
\end{proposition}

\begin{theorem}[Replay Equivalence]
The Spherepop compiler satisfies history equivalence for all programs expressible
in the Spherepop core calculus.
\end{theorem}

\begin{proof}
By structural induction on program terms. For each primitive operation (Refuse,
Bind, Collapse, Pop), the compiler emits the corresponding $\mathsf{IrEvent}$
variant, and the VM step performs exactly the same World mutation as the
interpreter's eval case. The history appended is therefore identical. For $\mathsf{Seq}$:
each sub-term is compiled in order, and since histories compose by concatenation,
the composed history of the compiled form equals that of the interpreted form by
induction.
\end{proof}

\begin{theorem}[Compiler Full Faithfulness]
\[
I(p_1).\mathsf{history} = I(p_2).\mathsf{history}
\iff
V(C(p_1)).\mathsf{history} = V(C(p_2)).\mathsf{history}.
\]
\end{theorem}

\begin{proof}
($\Rightarrow$) By Replay Equivalence applied twice. ($\Leftarrow$) Similarly.
The biconditional follows from the symmetry of history equivalence.
\end{proof}

\chapter{Garbage Collection as Admissibility-Guided Collapse}

\section{The Standard Problem}

In conventional languages, garbage collection identifies and reclaims memory that
is no longer reachable from the program's roots. The criterion is purely
structural: an object is garbage if no live reference points to it.

In Spherepop, the analogous question is: which portions of the history are no
longer needed to support any reachable continuation? The answer is richer because
it depends not just on structural reachability but on admissibility.

\section{History Segments and Liveness}

\begin{definition}[History Segment Liveness]
A history segment $H[i:j] = [e_i, e_{i+1}, \ldots, e_j]$ is \emph{live at
time $t$} if there exists an agent $A$ and a continuation $c$ such that:
\begin{enumerate}
\item $c$ is reachable at time $t$: $P_A(c, t) > 0$, and
\item $c$ depends on $H[i:j]$: $P_A(c, t \mid \neg H[i:j]) < P_A(c, t)$.
\end{enumerate}
A segment is \emph{dead} if it is not live.
\end{definition}

\begin{definition}[Admissibility-Guided GC]
\emph{Admissibility-guided garbage collection} removes dead history segments
while preserving the admissibility certificates needed for future collapses.
Formally, let $\calA_{\mathrm{live}}(H)$ be the set of future continuations
still reachable. GC produces $H' \subseteq H$ (as a subsequence) satisfying:
\[
\calA(H') = \calA(H) \cap \{x : \exists c \in \calA_{\mathrm{live}}(H),\, x \text{ needed by } c\}.
\]
\end{definition}

\section{GC as Collapse}

\begin{theorem}[GC as Collapse Composition]
Admissibility-guided GC is equivalent to applying a collapse rule $c_{\mathrm{GC}}$
that maps $H$ to the subsequence of live events:
\[
c_{\mathrm{GC}}(H) = [e \in H : e \text{ is live in } H].
\]
$c_{\mathrm{GC}}$ is a valid collapse rule when liveness is stable under
extension: if $e$ is live in $H$, it is live in $H \mathbin{++} H'$ for all $H'$.
\end{theorem}

\begin{proof}
Liveness is defined in terms of reachable continuations. Since histories grow
monotonically and Pop operations only add to the committed history, a past
commitment that enables a reachable continuation remains enabling under history
extension. Therefore $c_{\mathrm{GC}}$ respects composition: the GC of a
concatenated history equals the concatenation of the GC of the parts.
\end{proof}

\begin{corollary}[GC Preserves Admissibility]
If $c_{\mathrm{GC}}(H) = H'$, then $\calA(H') \supseteq \calA_{\mathrm{live}}(H)$.
\end{corollary}

\begin{proof}
GC removes only dead events. A dead Refuse event is one that marks a symbol
$x$ as inadmissible when $x$ is no longer needed by any reachable continuation.
Removing it restores $x$ to the admissibility set. Since no live continuation
needs $x$, this does not affect any future collapse certificate.
\end{proof}

\section{Three GC Strategies}

Different collapse rules for GC yield different strategies with different
cost-benefit profiles.

\begin{definition}[Mark-Scan GC]
\emph{Mark-Scan GC} applies the identity rule $c_I$ to the subhistory of live
events: it retains the full sequence of live events in order. This is the most
informative strategy and preserves all historical structure needed for replay.
\end{definition}

\begin{definition}[Quotient GC]
\emph{Quotient GC} applies a coarser rule $c_{LW}$ to the history before
collecting: it retains only the last commitment to each symbol. This is equivalent
to state-based GC with a history log and loses intermediate state information.
\end{definition}

\begin{definition}[Certificate-Only GC]
\emph{Certificate-Only GC} retains only Refuse events (inadmissibility
certificates) and Collapse events (observation certificates), discarding all Pop
and Bind events whose certificates have been sealed. This minimizes memory use
while preserving the admissibility structure needed for future type checking.
\end{definition}

\begin{theorem}[GC Strategy Tradeoff]
Let $M(c_{\mathrm{GC}})$ denote memory usage and $R(c_{\mathrm{GC}})$ denote
recovery power (the set of continuations that can be reconstructed after GC).
Then:
\[
M(c_I) \geq M(c_{LW}) \geq M(c_{\mathrm{cert}})
\quad \text{and} \quad
R(c_I) \geq R(c_{LW}) \geq R(c_{\mathrm{cert}}).
\]
\end{theorem}

\begin{proof}
The identity rule retains the most events and therefore uses the most memory but
also preserves the most recovery information. Each coarsening reduces both memory
usage and recovery power by identifying more history segments as equivalent.
\end{proof}

\section{Admissibility Certificates and GC Safety}

A key constraint on GC in Spherepop is that it must not invalidate pending type
judgements. If a collapse term $\mathsf{collapse}(t, c)$ is waiting on the
admissibility certificate $\Gamma \vdash t : \Adm(T)$, GC must retain the history
events that support that certificate.

\begin{definition}[GC Safety]
A GC operation is \emph{safe} if for every pending collapse judgement
$\Gamma \vdash \mathsf{collapse}(t, c) : \Col(T, c)$, the history events needed
to support $\Gamma \vdash t : \Adm(T)$ are retained after GC.
\end{definition}

\begin{theorem}[Safety of Certificate-Only GC]
Certificate-Only GC is safe: it retains all Refuse events that contribute to
any pending admissibility certificate.
\end{theorem}

\begin{proof}
The admissibility type $\Adm(T)$ is established by the sequence of Pop and Refuse
events that produced $t$. Certificate-Only GC retains all Refuse events
(inadmissibility certificates) unconditionally. The Pop events can be inferred
from the Collapse certificates that seal their results. Therefore all information
needed for pending collapse judgements is preserved.
\end{proof}

\chapter{Error Correction as Certified Refusal}

\section{The Standard Approach and Its Limits}

In conventional languages, error handling is an afterthought: exceptions are
thrown when something goes wrong, and error handling code is written separately
from the main computation. This approach has three problems. First, it treats
errors as exceptional when they are often structurally inevitable. Second, it
provides no information about why the error occurred. Third, it makes error
recovery expensive because the history of the computation is not available.

In Spherepop, errors are not exceptional. Every Refuse event is a documented
non-traversal of a possible path. Error handling is not separate from the main
computation; it is part of the computation's history.

\section{Error Correction Codes in History Space}

\begin{definition}[History Redundancy]
A history $H$ has \emph{redundancy factor} $k$ if there exist $k$ disjoint
subsequences $H_1, \ldots, H_k \subseteq H$ (as event subsequences) such that
each $H_i$ alone is sufficient to reconstruct the full admissibility certificate
for the terminal collapse:
\[
\calA(H_i) \supseteq \calA_{\mathrm{needed}}(H) \quad \text{for each } i.
\]
\end{definition}

\begin{definition}[Erasure-Correcting History]
An \emph{erasure-correcting history} with parameters $(n, k)$ encodes $k$
independent certificate paths through $n$ events, such that any $k$ of the $n$
events are sufficient to recover the full admissibility structure.
\end{definition}

\begin{theorem}[Error Correction via Redundant Refusal]
If a history $H$ with redundancy factor $k \geq 2$ has $m < k/2$ events
corrupted or erased, the correct admissibility structure can be recovered from
the remaining events.
\end{theorem}

\begin{proof}
With redundancy $k$ and fewer than $k/2$ erasures, at least $k - k/2 > k/2$
redundant paths remain intact. Majority vote over the surviving paths recovers
the correct admissibility judgement.
\end{proof}

\section{Typed Error Propagation}

\begin{definition}[Error Type]
An \emph{error type} is a type of the form $\Rfsd(T, r)$ where $r$ is a
proof-carrying reason. Error types carry the evidence of inadmissibility as
first-class type-level information.
\end{definition}

\begin{definition}[Error Propagation Rule]
\[
\frac{\Gamma \vdash t : \Rfsd(T_1, r) \quad \Gamma, x : T_1 \vdash b : T_2}{\Gamma \vdash \mathsf{handle}(t, x.b) : T_2 \vee \Ref(T_2, r)}
\quad [\mathrm{T\text{-}HANDLE}]
\]
The handler $b$ may produce a value of $T_2$ or propagate a refusal. The
propagated refusal carries the original reason $r$, preserving the full provenance
chain.
\end{definition}

\begin{theorem}[Error Preservation]
If $\Gamma \vdash t : \Rfsd(T, r)$ and there is no handler in scope for $r$,
then $r$ propagates to the nearest enclosing context that can handle
$\Ref(\cdot, r)$.
\end{theorem}

\begin{proof}
By the typing rules, $\Rfsd(T, r)$ types only propagate upward through
$[\mathrm{T\text{-}HANDLE}]$ if the handler does not fully resolve them. Since
every unresolved $\Rfsd(T, r)$ in a well-typed program must eventually reach a
handler or become the return type of the program, propagation terminates at the
program boundary with a documented reason.
\end{proof}

\section{Recovery as Certified Continuation}

\begin{definition}[Recovery Certificate]
A \emph{recovery certificate} for an error $\Rfsd(T, r)$ at world state $W$ is a
term $\mathsf{recover} : \Rfsd(T, r) \to \Adm(T')$ that transforms a documented
inadmissibility into an admissible continuation of type $T'$.
\end{definition}

\begin{theorem}[Soundness of Recovery]
If $\mathsf{recover} : \Rfsd(T, r) \to \Adm(T')$ is a recovery certificate and
$\Gamma \vdash t : \Rfsd(T, r)$, then $\Gamma \vdash \mathsf{recover}(t) : \Adm(T')$.
Furthermore, the history of $\mathsf{recover}(t)$ contains the full provenance
chain from the original refusal to the recovery, ensuring that the recovery is
not a silent erasure of the error.
\end{theorem}

\begin{proof}
By application of the recovery function type and $[\mathrm{T\text{-}POP}]$ applied
to the result. The history contains the original $\Refuse(x, r)$ event followed
by the recovery events. No history erasure is possible since histories grow
monotonically.
\end{proof}

\begin{remark}
Recovery in Spherepop is the computational analogue of repair theory: it restores
access to a continuation after documented inadmissibility, but unlike silent error
suppression, it preserves the full record of what went wrong and how it was
addressed. The recovery certificate is a note---specifically a Repair Note---in
the formal sense of the note taxonomy.
\end{remark}

%% ============================================================
%% PART VII: DEEP STRUCTURE
%% ============================================================
\part{Deep Mathematical Structure}

\chapter{The Observation Limit}

\section{The Reviewer's Challenge}

Spherepop does not abolish the state illusion. It makes the quotient explicit,
parameterised, and reversible whenever the underlying history is retained.

A careful reader will notice a tension. The document claims to resolve the state
illusion---the loss of historical information caused by projecting history to
state. Yet every output of a Spherepop program is a collapse: a value derived by
applying a rule to a history. The programmer observes the collapsed value, not
the history. Is Spherepop not simply relocating the state illusion?

This chapter addresses that challenge. We prove the No Direct Observation Theorem,
which shows the challenge identifies a mathematical necessity, and then distinguish
precisely between state erasure (which Spherepop eliminates) and state illusion
(which no computational system can eliminate).

\section{The No Direct Observation Theorem}

\begin{definition}[History-Faithful Observation]
An observation map $\phi : \calH \to O$ is \emph{history-faithful} if it is
injective: $\phi(H_1) = \phi(H_2) \Rightarrow H_1 = H_2$.
\end{definition}

\begin{theorem}[No Direct Observation]
Let $O$ be any set with $|O| < |\calH|$. Then every observation map
$\phi : \calH \to O$ is non-injective: there exist distinct $H_1 \neq H_2$ with
$\phi(H_1) = \phi(H_2)$. Consequently, no computational system with a finite
observable space can directly verify the history it constructed.
\end{theorem}

\begin{proof}
Since $\calH = \bigcup_{n \geq 0} E^n$ is countably infinite and $|O| < |\calH|$,
by the pigeonhole principle $\phi$ cannot be injective. Therefore there exist
$H_1 \neq H_2$ with $\phi(H_1) = \phi(H_2)$.
\end{proof}

\begin{remark}
Theorem 7.1 is a fundamental limit. It says: every observation is a quotient.
This is not a contingent limitation that better engineering can overcome. It is a
cardinality constraint. A system that made its observation map injective---that
recorded every distinction in full history space---would itself be a history.
\end{remark}

\section{Erasure versus Illusion}

\begin{definition}[State Erasure]
A computational system suffers \emph{state erasure} if it irrecoverably discards
the history $H$ after computing $\sigma(H)$. After execution, $H$ is unavailable
and only $\sigma(H)$ remains.
\end{definition}

\begin{definition}[State Illusion]
A computational system presents a \emph{state illusion} to its programmers if
they can only observe $\sigma(H)$ and not $H$ directly. By Theorem 7.1, some
state illusion is unavoidable for any non-trivial observable space.
\end{definition}

\begin{proposition}[What Spherepop Eliminates]
Spherepop eliminates state erasure: the history $H$ is retained in
\texttt{world.history} and is available to any observer. Spherepop cannot and
does not eliminate the state illusion: the programmer's primary observable is
$c(H)$, not $H$, so distinct histories that agree under $c$ remain
indistinguishable from the perspective of that observation.

The correct claim is therefore: Spherepop does not abolish the state illusion.
It makes the quotient explicit, parameterised, and history-preserving.
\end{proposition}

\section{Rule Refinement as Disambiguation}

\begin{definition}[Rule Refinement Order]
Collapse rule $c_1$ is \emph{finer than} $c_2$, written $c_1 \preceq c_2$, if
$H_1 \sim_{c_1} H_2 \Rightarrow H_1 \sim_{c_2} H_2$. The Identity rule $c_I$ is
the finest rule; the trivial rule $c_\top$ mapping all histories to a single point
is the coarsest.
\end{definition}

\begin{proposition}[Disambiguation by Refinement]
Let $H_1 \sim_c H_2$ (histories indistinguishable under rule $c$) and $H_1 \neq H_2$.
Then there exists a finer rule $c' \preceq c$ such that $c'(H_1) \neq c'(H_2)$.
In the limit, $c_I$ always disambiguates since $c_I(H) = H$ exactly.
\end{proposition}

\begin{proof}
Take $c' = c_I$: since $H_1 \neq H_2$ and $c_I(H) = H$, we have
$c_I(H_1) = H_1 \neq H_2 = c_I(H_2)$.
\end{proof}

\chapter{Category Theory as Engine}

\section{The History-Observation Adjunction}

\begin{definition}[History Construction Functor]
Let $\mathbf{Terms}$ be the category whose objects are Spherepop terms and whose
morphisms are reduction sequences. Define $G : \mathbf{Terms} \to \Hist$ by
$G(t) = I(t).\mathsf{history}$.
\end{definition}

\begin{definition}[Observation Functor]
For a fixed collapse rule $c$, define $\mathcal{O}_c : \Hist \to \Obs_c$ by
$\mathcal{O}_c(H) = c(H) = F_c(H)$.
\end{definition}

\begin{theorem}[History-Observation Adjunction]
For the Identity rule $c_I$, there is a natural bijection
\[
\Hist(G(t), H) \cong \mathbf{Terms}(t, G^{-1}(H))
\]
whenever $G^{-1}(H)$ is defined. This is the categorical statement of Replay
Equivalence: the interpreter and compiler produce histories that are morphisms in
$\Hist$.
\end{theorem}

\section{Collapse Rules Form a Lattice}

\begin{proposition}[Observational Lattice]
The set of collapse rules ordered by refinement $c_1 \preceq c_2$ forms a complete
lattice. The meet of $\{c_i\}$ is the finest rule at least as fine as all; the
join is the rule generated by the union of equivalence relations.
\end{proposition}

\begin{proof}
The meet and join are defined by the standard lattice operations on equivalence
relations. Completeness follows because arbitrary intersections and generated
equivalence relations of equivalence relations are again equivalence relations.
\end{proof}

\section{Sheaf-Theoretic Semantics}

\begin{definition}[Observational Site]
Fix collapse rules $\{c_i\}$ and observable states $\{o_i\}$. The
\emph{observational site} $\mathcal{O}$ is the category whose objects are
observational regions $\{c_i^{-1}(o_i)\}$ and whose morphisms are inclusions
$c_j^{-1}(o_j) \hookrightarrow c_i^{-1}(o_i)$ when the former is contained in
the latter.
\end{definition}

\begin{proposition}[$\calH$ is a Sheaf]
The history presheaf $\calH$ on $\mathcal{O}$---defined by $\calH(U) = \{H \in \calH : H \in U\}$---satisfies the sheaf condition: compatible local histories glue
uniquely to global histories.
\end{proposition}

\begin{proof}
Compatibility means event sequences agree on their shared prefixes. Two compatible
event sequences agreeing on all overlaps determine a unique concatenation.
\end{proof}

\begin{proposition}[State Illusion as Non-Injectivity of Sections]
A collapse rule $c$ exhibits state illusion if and only if the observable sheaf
$\mathcal{O}_c$ fails to determine a unique global section of $\calH$. There
exist distinct $H_1, H_2 \in \calH(\calH)$ with the same section in
$\mathcal{O}_c(\calH)$.
\end{proposition}

The sheaf structure ties the local (what one observer sees under one rule) to the
global (the full history that all rules project from). Sheaf cohomology measures
the obstruction to lifting local observations to global histories---which is
precisely the observational entropy.

\chapter{CLIO Projections and Active Geodesic Inference}

\section{Collapse Rules as CLIO Projections}

\begin{definition}[CLIO Projection]
A CLIO projection $\pi : X \to M$ maps a representation space $X$ to a manifold
$M$ of observational strata. In the Spherepop setting: $X = \calH$ (history
space), $M = O_c$ (observable state space under rule $c$), and $\pi = F_c$
(the collapse functor).
\end{definition}

\begin{definition}[Observational Entropy]
For a collapse rule $c$ and observable state $o \in O_c$, the
\emph{observational entropy} is
\[
S_c(o) = \log |c^{-1}(o)|.
\]
\end{definition}

\begin{theorem}[Entropy Monotonicity Under Rule Refinement]
If $c_1 \preceq c_2$ ($c_1$ finer), then for every $o \in O_{c_1}$:
\[
S_{c_1}(o) \leq S_{c_2}(c_2(c_1^{-1}(o))).
\]
Finer rules have smaller or equal fibre sizes.
\end{theorem}

\begin{proof}
Since $c_1 \preceq c_2$, every $\sim_{c_1}$-class is contained in a
$\sim_{c_2}$-class. Therefore $|c_1^{-1}(o)| \leq |c_2^{-1}(c_2(H))|$ for
any $H \in c_1^{-1}(o)$.
\end{proof}

\begin{remark}
Observational entropy $S_c(o) = \log |c^{-1}(o)|$ is exactly the CLIO quantity
$S_\pi(m) = \log \mathrm{Vol}(\pi^{-1}(m))$ in the discrete Spherepop setting.
Spherepop is a discrete computational instance of the CLIO projection program.
\end{remark}

\section{The Santoro--Waghmare--Panaretos Unification}

\begin{theorem}[Representation Changes Topology]
Let $\rho_1, \rho_2 : X \to Y$ be two representation maps with $\tau_{\rho_1} \neq \tau_{\rho_2}$. Then there exist $x_1, x_2 \in X$ such that $x_1 \sim_{\rho_1} x_2$
but $x_1 \not\sim_{\rho_2} x_2$, or vice versa. The two representations disagree
on whether $x_1$ and $x_2$ are the same object.
\end{theorem}

\begin{proof}
If $\tau_{\rho_1} \neq \tau_{\rho_2}$, there exists an open set
$U \in \tau_{\rho_1} \setminus \tau_{\rho_2}$. Then $U = \rho_1^{-1}(V)$ for
some $V \subseteq Y$, but $U \neq \rho_2^{-1}(W)$ for any $W$. Therefore there
exist $x_1, x_2$ with $\rho_1(x_1) = \rho_1(x_2)$ but $\rho_2(x_1) \neq \rho_2(x_2)$.
\end{proof}

\begin{remark}
This theorem is the general form of the observation that representations
reorganise topology. It applies uniformly to:
(1) the Santoro--Waghmare--Panaretos kernel embedding, which converts metrically
close distributions into topologically singular Gaussian measures; and
(2) Spherepop collapse rules, where refinement moves history pairs between
equivalence classes.

In the kernel case, the topological change is forced by the Feldman--H\'ajek
theorem. In the collapse case, the topological change is chosen: the programmer
selects the rule, thereby choosing which distinctions are topologically visible.
Both are instances of the same structural operation.
\end{remark}

\section{Active Geodesic Inference}

The Spherepop framework connects to a broader theory of intelligent systems as
trajectory-maintaining structures.

\begin{definition}[Active Geodesic Inference]
An intelligent system performs \emph{active geodesic inference} if it actively
shapes its configuration space through energetic and entropic constraints, forming
low-action trajectories through possibility space rather than following
predetermined state transitions.
\end{definition}

\begin{proposition}[Spherepop as Active Geodesic Substrate]
Spherepop's irreversible, scope-based semantics enforce history-sensitivity and
prevent incoherent superposition by construction. Pop operations commit to
specific trajectories; Refuse operations document inadmissible branches; the
history preserves the full record of path selection. Spherepop is therefore a
natural substrate for implementing active geodesic inference: the history is the
geodesic; the admissibility set is the constraint field; collapse is the
observation that seals a trajectory as complete.
\end{proposition}

\begin{definition}[Formal Intelligence]
Intelligence is the capacity of a system to remain within a family of dynamically
admissible, low-action histories by actively reshaping its configuration space's
geometry. This definition extends across scales, encompassing learning within a
lifetime, evolution across generations, and reasoning within an episode.
\end{definition}

\begin{theorem}[Intelligence as Note Infrastructure]
An intelligent system with formal intelligence in the above sense is necessarily
a note-producing system: it constructs and maintains Capture, Prospective, and
Repair notes as part of the mechanism by which it reshapes its configuration space.
Notes are the mechanism by which intelligent systems install future reachability
from present positions.
\end{theorem}

\begin{proof}
A system that maintains low-action histories across time must preserve information
about its past trajectories (Capture Notes), commit to future trajectories
(Prospective Notes), and recover from trajectory interruptions (Repair Notes).
Each of these is a note class in the formal taxonomy. Therefore any formally
intelligent system produces notes as a structural necessity, not a contingent
feature.
\end{proof}

%% ============================================================
%% PART VIII: CIVILIZATION
%% ============================================================
\part{Civilization as Note Infrastructure}

\chapter{MEM\textbar8 and the Cognitive Substrate}

\section{Memory as Event Structure}

The MEM\textbar8 architecture treats memory not as isolated symbol storage but as
structured pattern access where retrieval depends on preserving relationships,
locality, and accessibility. Memory is not a database; it is a cognitive substrate
whose operations are structurally similar to Spherepop's history operations.

\begin{proposition}[MEM\textbar8 as Spherepop Instance]
The MEM\textbar8 architecture is a biological instance of the Spherepop framework:
\begin{center}
\begin{tabular}{ll}
\toprule
MEM\textbar8 Layer & Spherepop Note Class \\
\midrule
Episodic layer & Capture notes \\
Procedural layer & Repair notes \\
Semantic layer & Coordination notes \\
Indexical layer & Collapse notes \\
Theoretical layer & Generative notes \\
\bottomrule
\end{tabular}
\end{center}
\end{proposition}

\begin{remark}
A notebook is a miniature MEM\textbar8. A Git repository is a miniature
MEM\textbar8. An archive is a miniature MEM\textbar8. Each stores structured
trajectories organized to make reconstruction possible. The difference between
a good archive and a poor one is not quantity of material stored but quality
of access structure: how well the collapse notes and coordination notes that
organize the collection preserve the reachability of trajectories within it.
\end{remark}

\chapter{Repair Theory and Civilizational Collapse}

\section{Repair as Pre-emptive Refusal}

Repair Theory holds that repair is the restoration of access to a trajectory
that has been interrupted. The connection to Spherepop is not metaphorical but
structural: repair is the execution of a Repair Note that was created in
anticipation of the failure.

\begin{theorem}[Note Taxonomy as Failure Taxonomy]
The seven note classes correspond bijectively to seven classes of anticipated
trajectory failure:
\begin{center}
\begin{tabular}{ll}
\toprule
Note Class & Anticipated Failure \\
\midrule
Capture & Experiential collapse \\
Prospective & Intention-execution gap \\
Repair & Procedural interruption \\
Coordination & Coordination failure \\
Refusal & Distinction collapse \\
Generative & Inferential poverty \\
Collapse & Navigational failure \\
\bottomrule
\end{tabular}
\end{center}
\end{theorem}

\section{Civilizational Scale}

\begin{definition}[Civilizational Note Infrastructure]
Let civilization $Z$ possess note infrastructure $\calN = \{N_1, \ldots, N_k\}$.
The \emph{collective reachability} is
\[
R(Z) = \bigcup_{i=1}^{k} \calC(N_i),
\]
where $\calC(N_i)$ is the set of continuations preserved or enabled by $N_i$.
\end{definition}

\begin{theorem}[Civilizational Collapse]
The destruction of note infrastructure $\calN$ decreases collective reachability
even when underlying truths remain unchanged:
\[
R(Z \setminus \calN) < R(Z).
\]
A civilization that loses its note infrastructure does not become ignorant. It
becomes unable to reach its own knowledge.
\end{theorem}

\begin{proof}
By definition, $R(Z)$ includes all continuations preserved by elements of $\calN$.
If $\calN$ is destroyed, no element contributes to collective reachability. Since
some continuations are accessible only through $\calN$, it follows that
$R(Z \setminus \calN) < R(Z)$. The underlying facts are not destroyed, but the
paths that made them accessible are.
\end{proof}

\begin{center}
\begin{tabular}{lll}
\toprule
Technology & Reachability Horizon & Dominant Note Classes \\
\midrule
Speech & Minutes to hours & Coordination, Capture \\
Story & Decades & Capture, Generative \\
Writing & Centuries & All classes \\
Printing & Populations & Coordination, Refusal \\
Digital storage & Replication & All classes, at scale \\
Version control & Modification history & Repair, Capture \\
Network & Distribution & Coordination, Generative \\
\bottomrule
\end{tabular}
\end{center}

Each note technology expands the reachability horizon by addressing failure modes
the preceding technology could not prevent. Writing prevents the death of
witnesses. Printing prevents the destruction of individual copies. Version control
prevents the loss of developmental history through modification. The history of
civilization is the history of expanding note infrastructure.

%% ============================================================
%% PART IX: OPEN PROBLEMS
%% ============================================================

%% ============================================================
%% PART IX: RECONSTRUCTION, REFUSAL, AND REPAIR
%% ============================================================
\part{Reconstruction, Refusal, and Repair}

\chapter{Observational Entropy and Reconstruction Cost}

\section{From Observation to Reconstruction}

The preceding chapters established that every collapse rule
$c : \calH \to O_c$ induces an observational entropy
$S_c(o) = \log |c^{-1}(o)|$.
This quantity measures the size of the fibre associated with an
observation. A collapse with small observational entropy preserves many
distinctions between histories. A collapse with large observational entropy
identifies many histories and destroys information. The natural next question is
not merely how many histories are hidden, but how difficult it is to reconstruct
the correct one.

\begin{definition}[Reconstruction Cost]
Let $o \in O_c$. The \emph{reconstruction cost} of $o$ is
\[
\calR(o) = \mathbb{E}\bigl[C(H) \mid H \in c^{-1}(o)\bigr]
\]
where $C(H)$ denotes the computational effort required to identify history $H$
from among the candidates in the fibre.
\end{definition}

\begin{definition}[Uniform Reconstruction]
A reconstruction procedure is \emph{uniform} if every history in $c^{-1}(o)$
is considered equiprobable given only the observation $o$.
\end{definition}

\begin{theorem}[Entropy Lower Bound on Reconstruction]
For any uniform reconstruction procedure,
\[
\calR(o) \geq k \cdot S_c(o)
\]
for some positive constant $k$ depending only on the cost model $C$.
\end{theorem}

\begin{proof}
The fibre $c^{-1}(o)$ contains $N = |c^{-1}(o)|$ candidate histories. Any
reconstruction procedure must distinguish among these alternatives. A decision
tree distinguishing $N$ equally probable alternatives requires at least $\log_2 N$
binary decisions. Each decision has a positive cost bounded below by $k > 0$.
Since $S_c(o) = \log N$, the expected reconstruction effort satisfies
$\calR(o) \geq k \log N = k \cdot S_c(o)$.
\end{proof}

\begin{corollary}
Observational entropy is a lower bound on reconstruction difficulty. Notes that
reduce observational entropy are therefore notes that reduce reconstruction cost.
\end{corollary}

\begin{remark}
This theorem explains precisely why notes have value in terms the continuation
theory can measure. A note $N$ reduces $|c^{-1}(o)|$ by introducing additional
distinctions that collapse cannot erase. Reducing the fibre size reduces $S_c(o)$
and hence reduces $\calR(o)$. The note acts as a compression-resistant continuation
preserver: it encodes information that the observational collapse would otherwise
destroy, lowering the expected work required to recover the original history.
\end{remark}

\section{Continuation Value and Note Leverage}

\begin{definition}[Continuation Value of an Observation]
The \emph{continuation value} of observation $o$ is
\[
V(o) = \frac{1}{1 + \calR(o)}.
\]
\end{definition}

Observations with large fibres have low continuation value because the
original history is expensive to recover. Observations with small fibres have
high continuation value. The note taxonomy from Part~V can now be given a
quantitative characterization: a note $N$ is valuable to the extent that it
increases $V(o)$ for observations $o$ it affects, which is equivalent to
decreasing $\calR(o)$, which is equivalent to decreasing $S_c(o)$.

\begin{theorem}[Note Leverage as Entropy Reduction]
The reachability leverage $\Lambda(N)$ of a note (Definition 2.3) is bounded
below by the entropy reduction it induces:
\[
\Lambda(N) \geq \frac{k}{C_{\mathrm{create}}(N)} \sum_{o} [S_c(o) - S_c(o \mid N)],
\]
where $S_c(o \mid N)$ denotes the entropy of observation $o$ when the note $N$
is available.
\end{theorem}

\begin{proof}
By the Entropy Lower Bound theorem, entropy reduction lower-bounds reconstruction
cost reduction. Summing over affected observations and dividing by creation cost
gives the claimed bound on leverage.
\end{proof}

\section{The Reconstruction Complexity Hierarchy}

Different collapse rules induce different reconstruction complexity classes.

\begin{definition}[Tractable Reconstruction]
A collapse rule $c$ has \emph{tractable reconstruction} if $\calR(o)$ is bounded
by a polynomial in the size of the history alphabet for all $o \in O_c$.
\end{definition}

\begin{proposition}
The Identity rule $c_I$ has reconstruction complexity $O(1)$: given the
observation (which equals the history), reconstruction is immediate. The
LastWrite rule $c_{LW}$ has reconstruction complexity exponential in the
number of symbols, since the fibre contains all histories consistent with the
last-write state. The Accumulate rule $c_{\mathrm{acc}}$ has reconstruction
complexity equal to the multinomial coefficient $\binom{|\text{events}|}{n_1, \ldots, n_k}$
where $n_i$ is the count for symbol $i$.
\end{proposition}

\chapter{The Geometry of Refusal}

\section{Refusal as Geometric Excision}

Refusal was introduced in Part~II as a documented declaration of inadmissibility.
A deeper interpretation emerges when admissibility is viewed as a region of
possibility space. The admissibility set $\calA(H)$ is not merely a set of
available symbols; it is a geometric region in possibility space, and each refusal
excises a portion of that region.

\begin{definition}[Admissibility Region]
The \emph{admissibility region} after history $H$ is
$\calA(H) = \Omega_0 \setminus \{x \mid \exists r.\, \Refuse(x,r) \in H\}$.
\end{definition}

\begin{definition}[Refusal Region]
A refusal reason $r$ determines a \emph{forbidden subset} $U_r \subseteq \Omega_0$.
A refusal $\Refuse(x, r)$ excises $x$ from the admissibility region under
the constraint $r$.
\end{definition}

\begin{definition}[Refusal Action]
Applying $\Refuse(x, r)$ produces the new admissibility region
\[
\calA(H') = \calA(H) \setminus \{x\}.
\]
When $r$ is a constraint that applies to a family of symbols
$U_r = \{x : r(x) = \mathsf{true}\}$, the excision generalizes to
$\calA(H') = \calA(H) \setminus U_r$.
\end{definition}

Refusal is therefore a geometric excision operator: it removes regions from
possibility space. Multiple refusals compose to produce an increasingly constrained
admissibility region, a process that can be analyzed in terms of the topology of
the resulting space.

\section{Refusal Monotonicity}

\begin{theorem}[Refusal Monotonicity]
For every refusal event, $\calA(H') \subseteq \calA(H)$.
\end{theorem}

\begin{proof}
By construction, $\calA(H') = \calA(H) \setminus U_r$ for some $U_r \subseteq \Omega_0$.
Removing a subset cannot enlarge the region. Hence $\calA(H') \subseteq \calA(H)$.
\end{proof}

\begin{corollary}[Irreversibility of Refusal]
Once a symbol $x$ is refused, it cannot be re-admitted without replacing the
history. Since histories grow monotonically, refusal is irreversible within a
computation.
\end{corollary}

\section{Refusal Curvature}

\begin{definition}[Local Admissibility Curvature]
For a point $x$ in possibility space with neighbourhood $N(x)$,
the \emph{local admissibility curvature} is
\[
\kappa_A(x) = \frac{|N(x) \cap \calA|}{|N(x)|}.
\]
A curvature of $1$ indicates full local admissibility; a curvature of $0$
indicates a refusal singularity.
\end{definition}

\begin{definition}[Refusal Singularity]
A point $x$ is a \emph{refusal singularity} if $\kappa_A(x) = 0$: every
trajectory through $x$ is inadmissible.
\end{definition}

\begin{proposition}[Curvature Monotonicity Under Refusal]
Repeated refusals create regions of non-increasing admissibility curvature:
if $H' = H \mathbin{++} [\Refuse(x, r)]$ then $\kappa_{A'}(y) \leq \kappa_A(y)$
for all $y \in U_r$.
\end{proposition}

\begin{proof}
Each refusal removes admissible trajectories from $U_r$. Therefore
$|N(y) \cap \calA'| \leq |N(y) \cap \calA|$ for $y \in U_r$, which implies
$\kappa_{A'}(y) \leq \kappa_A(y)$.
\end{proof}

\begin{definition}[Admissibility Field]
The function $\kappa_A : \Omega_0 \to [0,1]$ is the \emph{admissibility field}
of the computation. Its level sets $\{x : \kappa_A(x) \geq \lambda\}$ for
threshold $\lambda$ define the $\lambda$-admissible regions of possibility space.
\end{definition}

\begin{remark}
The admissibility field $\kappa_A$ is the discrete computational analogue of
the smooth admissibility field in the RSVP framework. The xylomorphic criterion
$\lambda < 1$ from RSVP corresponds to $\kappa_A(x) < 1$ in the discrete
setting: a trajectory is admissible at level $\lambda$ when its local admissibility
curvature exceeds $\lambda$. The RSVP continuum limit of the Spherepop admissibility
field is therefore obtained by taking the lattice spacing to zero in the discrete
geometry.
\end{remark}

\section{The Galois Connection Between History Growth and Admissibility Contraction}

\begin{theorem}[History--Admissibility Galois Connection]
There is a Galois connection between the lattice of history extensions and the
lattice of admissibility regions, ordered by the relation:
\[
H \leq H' \iff H \text{ is a prefix of } H'
\]
on histories and $\calA \leq \calA'$ iff $\calA \subseteq \calA'$ on admissibility
regions. The Galois connection is given by:
\[
\Phi(H) = \calA(H) \quad \text{and} \quad \Psi(\calA) = \{H : \calA(H) \supseteq \calA\}.
\]
\end{theorem}

\begin{proof}
We verify the Galois condition: $H \leq H'$ implies $\Phi(H) \geq \Phi(H')$
(order reversal). By Refusal Monotonicity, if $H'$ extends $H$ by any sequence
of events, then $\calA(H') \subseteq \calA(H)$, that is $\Phi(H') \leq \Phi(H)$.
Conversely, $\calA \geq \calA'$ (i.e.\ $\calA \supseteq \calA'$) implies
$\Psi(\calA) \leq \Psi(\calA')$ (shorter histories needed to maintain the larger
admissibility set). The adjunction conditions $H \leq \Psi(\Phi(H))$ and
$\calA \leq \Phi(\Psi(\calA))$ follow from the definitions.
\end{proof}

\begin{remark}
The Galois connection formalizes the intuition that history growth and admissibility
contraction are dual processes. As a computation proceeds (history grows), the
space of admissible futures contracts. The Galois connection makes this duality
precise: every increase in historical commitment corresponds to a decrease in
future admissibility, and vice versa.
\end{remark}

\chapter{Notes as Externalized Adjoints}

\section{The Reconstruction Gap}

Collapse maps histories into observations, losing information. Notes partially
reverse this process by preserving information that collapse discards. The
question is: can this reversal be made precise in category-theoretic terms?

The answer is yes, and the result gives the note theory its deepest mathematical
formulation. Notes are partial right adjoints to collapse functors.

\section{The Note Reconstruction Functor}

\begin{definition}[Note Reconstruction Functor]
Given a collapse functor $\mathcal{O}_c : \Hist \to \Obs_c$, a
\emph{note reconstruction functor} is a functor
$\mathcal{N} : \Obs_c \to \Hist$
that assigns to each observable state $o$ a canonical history
$\mathcal{N}(o) \in c^{-1}(o)$.
\end{definition}

\begin{definition}[Adjoint Note]
A note is \emph{adjoint to collapse} if
\[
\Obs_c(\mathcal{O}_c(H), o) \cong \Hist(H, \mathcal{N}(o))
\]
naturally in $H$ and $o$.
\end{definition}

\begin{theorem}[Notes as Partial Right Adjoints]
\label{thm:notes-adjoints}
Every successful note acts as a partial right adjoint to a collapse operation,
in the sense that it provides a reconstruction map from observations back toward
histories on the subdomain where reconstruction succeeds.
\end{theorem}

\begin{proof}
Let $N$ be a note and $c$ a collapse rule such that $N$ preserves information
discarded by $c$ on some subdomain $D \subseteq O_c$. Define
$\mathcal{N}(o) = H_N(o)$ where $H_N(o)$ is the canonical history recovered
using note $N$ from observation $o$. For $o \in D$, the map
$\Hist(H, \mathcal{N}(o)) \to \Obs_c(\mathcal{O}_c(H), o)$ given by
$f \mapsto \mathcal{O}_c \circ f$ is an isomorphism: every history morphism
into $\mathcal{N}(o)$ projects to a unique observable morphism, and conversely,
since $N$ provides the additional data needed to lift observable morphisms to
history morphisms on $D$. This is the universal property of a right adjoint,
restricted to $D$.
\end{proof}

\begin{corollary}
The statement ``notes are anti-collapse artifacts'' is equivalent to saying that
notes approximate right adjoints to collapse functors. The richer the note's
continuation volume, the larger the subdomain $D$ on which the adjunction holds.
\end{corollary}

\section{Examples of Note Adjunctions}

The note-as-adjoint formulation unifies the taxonomy from Part~V.

A \emph{bookmark} is a right adjoint to the collapse of a URL into a browsing
context: it recovers the specific location from an otherwise unlabeled observation
of web content.

A \emph{mathematical proof} is a right adjoint to the collapse of a theorem
into a bare assertion: it recovers the inferential trajectory from the terminal
claim.

\emph{Source code} is a right adjoint to the collapse of executable behavior
into observable outputs: it recovers the computational trajectory from the
input-output specification.

A \emph{map} is a right adjoint to geographical collapse: it recovers spatial
structure from positional observations that contain no direction or distance
information.

In every case, the note $\mathcal{N}$ reconstructs distinctions erased by the
collapse $\mathcal{O}_c$, satisfying the adjunction on the subdomain where
reconstruction is possible.

\section{The Completeness Spectrum}

\begin{definition}[Note Completeness]
A note $N$ is \emph{complete} for collapse rule $c$ if $\mathcal{N}$ extends to a
global right adjoint: $\mathcal{O}_c \dashv \mathcal{N}$ on all of $\Obs_c$.
$N$ is \emph{partial} if the adjunction holds only on a proper subdomain.
\end{definition}

\begin{proposition}
No finite note is complete for a collapse rule that identifies infinitely many
distinct histories. Complete notes for finite history spaces exist and correspond
to observationally complete collapse rules (Proposition 19.3 of Part~VII).
\end{proposition}

\begin{proof}
If $c$ identifies infinitely many histories, then $|c^{-1}(o)| = \infty$ for
some $o$. A finite note contains finite information and can therefore eliminate
at most finitely many ambiguities per observation. Hence it cannot provide a
complete right adjoint on all of $\Obs_c$.
\end{proof}

\chapter{The Topology of Repair}

\section{Repair as a Morphism in the History Category}

Repair Theory holds that repair is the restoration of access to a trajectory
that has been interrupted. Within the Spherepop framework, repair can be
formalized as a special class of morphism in the history category.

\begin{definition}[Repair Morphism]
A \emph{repair morphism} from history $H$ to history $H'$ is a map
$\rho : H \to H'$ such that $c(H) = c(H')$ for some designated collapse rule
$c$: the observable behavior is preserved while the historical structure is altered.
\end{definition}

\begin{definition}[Repair Equivalence]
Two histories are \emph{repair equivalent}, written $H \sim_R H'$, if there
exists a collapse rule $c$ and a repair morphism $\rho : H \to H'$ such that
$c(H) = c(H')$.
\end{definition}

\begin{theorem}[Repair Equivalence Theorem]
Repair preserves observable reachability while modifying historical structure:
if $H \sim_R H'$ then every continuation observable under $c$ from $H$ is
also observable from $H'$, yet $H$ and $H' $ may differ as elements of $\calH$.
\end{theorem}

\begin{proof}
Since $c(H) = c(H')$, the two histories are in the same observational equivalence
class. Every continuation whose reachability depends only on the observable state
$c(H)$ is therefore equally reachable from $H'$. However, since $H \neq H'$ as
sequences of events, their internal structure differs. The repair has altered the
history while preserving its observable footprint.
\end{proof}

\section{Examples of Repair Classes}

The repair morphism formulation unifies several classes of real-world operation.

\emph{Software patches} replace a sequence of code events producing a buggy
observable behavior with a different sequence producing the correct observable
behavior. The patched and unpatched code share the same input-output specification
(the collapse rule is the black-box I/O function) but have different development
histories.

\emph{Biological healing} replaces a damaged tissue history with a regenerated
tissue history. The organism's observable function is preserved (the collapse rule
is the phenotypic function) while the internal molecular history changes.

\emph{Engineering maintenance} replaces a worn component history with a replacement
component history. The system's operational behavior is preserved (the collapse
rule is the performance specification) while the material history differs.

\emph{Institutional reform} replaces a history of governance decisions with a
revised sequence. The institution's external behavior and mandate are preserved
(the collapse rule is the institutional mandate) while the decision history changes.

All four are repair morphisms in the formal sense.

\section{Repair Cohomology}

Not every interrupted trajectory can be repaired. The obstruction to repair is
a topological invariant.

\begin{definition}[Repair Defect]
The \emph{repair defect} of histories $H$ and $H'$ is
$\Delta_R(H, H') = d_{\Hist}(H, H')$ where $d_{\Hist}$ measures the minimum
edit distance in the history category (minimum number of event replacements
needed to convert $H$ into a history with the same observable state as $H'$).
\end{definition}

\begin{definition}[Repair Cohomology Class]
When no repair morphism exists between $H$ and a target observable state $o$
(that is, $c^{-1}(o) \cap \mathsf{Reachable}(H) = \varnothing$), the
obstruction is a non-trivial element of the \emph{repair cohomology}
$H^1_R(\calH, c)$---the first cohomology group of the history space with
coefficients in the collapse sheaf.
\end{definition}

\begin{proposition}[Trivial Cohomology Implies Universal Repairability]
If $H^1_R(\calH, c) = 0$, then every history can be repaired to any
observationally equivalent target. The absence of cohomological obstructions
corresponds to a globally connected repair graph.
\end{proposition}

\begin{remark}
Repair cohomology provides a rigorous mathematical notion of irreparability.
A system with non-trivial repair cohomology has histories from which no
repair morphism can reach certain observable states, regardless of the effort
invested. In practice, this corresponds to: cryptographic irreversibility
(a hash cannot be repaired to recover the preimage), biological cell death
(a necrotic cell cannot be repaired to a living state), and data loss without
redundancy (a destroyed archive without backup admits no repair).
\end{remark}

\chapter{The Noun Fallacy Revisited}

\section{Objects as Stable Quotients}

Throughout this monograph we have argued that histories are primary and states
are derived. The present chapter makes this precise by giving a formal account
of how objects---the apparently primitive entities of everyday ontology---arise
from the collapse of histories.

\begin{definition}[Persistent Quotient]
An equivalence class $[H]_c$ under collapse rule $c$ is \emph{persistent} if
$c(H_t) = c(H_{t+\Delta})$ for all admissible extensions $H_{t+\Delta}$ of $H_t$
in a neighborhood of the current trajectory. That is, the observable state does
not change under small perturbations of the computation.
\end{definition}

\begin{definition}[Object]
An \emph{object} in the Spherepop ontology is a persistent quotient:
\[
\mathsf{Object} = [H]_c \quad \text{where } [H]_c \text{ is persistent.}
\]
\end{definition}

\section{Object Emergence}

\begin{theorem}[Object Emergence Theorem]
Stable observational equivalence classes emerge as objects to observers
who access histories only through collapse rules.
\end{theorem}

\begin{proof}
An observer accesses $\calH$ only through $c$. Repeated observation of the same
equivalence class produces observations $c(H_{t_1}) = c(H_{t_2}) = \cdots = o$.
The observer, receiving identical observations across time, infers an invariant
entity. This entity is $[H]_c$. The observer treats the persistent equivalence
class as an object because it is behaviorally indistinguishable from one: it
produces the same observable state under every observation the rule $c$ permits.
\end{proof}

\section{The Noun Fallacy}

\begin{definition}[Noun Fallacy]
The \emph{noun fallacy} is the error of treating a persistent quotient $[H]_c$
as an ontologically primitive object---as if $[H]_c$ existed independently of
the history $H$ and the collapse rule $c$ that produced it.
\end{definition}

\begin{corollary}
Every noun is a collapse note.
\end{corollary}

\begin{proof}
A noun compresses a family of trajectories into a single retrievable symbol.
It therefore performs the Collapse operation on the histories it subsumes. The
symbol is the collapsed observable; the suppressed histories are the fibres of the
collapse. A noun is precisely a Collapse Note in the sense of Chapter~13.
\end{proof}

\begin{example}
The word \emph{tree} is a collapse note. It collapses the family of all
developmental, ecological, metabolic, and evolutionary histories of tree-shaped
organisms into a single navigable token. The fibre $c_{\mathrm{tree}}^{-1}(\text{``tree''})$
is the set of all histories that produce observable states classifiable as trees.
The noun \emph{tree} makes this fibre navigable without requiring traversal of any
particular history within it.
\end{example}

\begin{example}
The word \emph{program} is a collapse note for computational history. It collapses
the set of all source texts, development decisions, debugging sessions, and
revision histories that produced a given executable into a single noun. The Spherepop
framework aims to reverse this collapse: to make the history that the noun
conceals explicitly accessible.
\end{example}

\section{Reality as Reachable History}

The preceding analysis permits a precise statement of the monograph's deepest claim.

\begin{theorem}[Reality as Reachability]
Under the Spherepop ontology:
\begin{enumerate}
\item Histories are primary: they are the fundamental objects of computation, cognition, and civilization.
\item Observations are quotients: they are equivalence classes of histories under collapse rules.
\item Objects are persistent quotients: stable equivalence classes that observers reify as things.
\item Notes preserve distinctions across quotients: they are artifacts that resist specific collapse operations.
\item Repair restores reachability after interruption: it constructs repair morphisms between interrupted and target histories.
\item Civilization constructs note infrastructure: it builds distributed systems for preserving continuations across generations.
\item Reality is a structured space of reachable histories, not a collection of primitive objects. Objects are what histories look like after sufficient collapse. 
\end{enumerate}
\end{theorem}

\begin{proof}
Items (1)--(6) follow from the formal definitions and theorems developed in
Parts II through VIII. Item (7) is the synthesis: given that objects are
persistent quotients (Definition 29.2) and quotients are derived from histories
(Proposition 2.2), objects depend ontologically on histories. The space of
objects is therefore a derived structure over the primary space of histories.
\end{proof}

%% ============================================================
%% PART X: COLLAPSE AS ORTHOGONAL PROJECTION
%% ============================================================
\part{Collapse as Orthogonal Projection}

\chapter{The Hilbert Space of Histories}

\section{Motivation}

Throughout the preceding chapters, collapse has been treated combinatorially:
a collapse rule $c : \calH \to O_c$ maps histories to equivalence classes, and
observational entropy counts fibre sizes. This combinatorial treatment is adequate
for finite history spaces, but it leaves the continuous and probabilistic dimensions
of collapse undeveloped.

The Hilbert space formulation of conditional expectation, developed in
probability theory through the work of Kolmogorov and later geometrized in
the $L^2$ framework, provides exactly the continuous analogue of Spherepop
collapse. The central thesis of this part is:

\begin{quote}
Every collapse rule introduced in previous chapters may be viewed as an orthogonal
projection operator. In finite settings, collapse appears combinatorial. In
probabilistic settings, collapse becomes orthogonal projection in a Hilbert space.
The two treatments are the same operation at different levels of abstraction.
\end{quote}

\section{The $L^2$ History Space}

\begin{definition}[$L^2$ History Space]
Let $(\Omega, \mathcal{F}, P)$ be a probability space where $\Omega$ is the
set of all histories, $\mathcal{F}$ is a $\sigma$-algebra over $\Omega$, and
$P$ is a probability measure over histories. The \emph{$L^2$ history space} is
\[
L^2(\mathcal{F}) = \{X : \Omega \to \mathbb{R} \mid \mathbb{E}[X^2] < \infty\}
\]
with inner product $\langle X, Y \rangle = \mathbb{E}[XY]$.
\end{definition}

\begin{remark}
Random variables $X \in L^2(\mathcal{F})$ are functions on the space of histories.
They measure observable quantities of histories---durations, counts of specific
event types, the time of the last Pop, the total weight of admissibility
certificates---in a way that respects the probability structure over histories.
\end{remark}

\begin{definition}[Observational Subspace]
Let $\mathcal{G} \subseteq \mathcal{F}$ be a sub-$\sigma$-algebra representing
an observational framework. The \emph{observational subspace} is
\[
L^2(\mathcal{G}) = \{X \in L^2(\mathcal{F}) \mid X \text{ is } \mathcal{G}\text{-measurable}\}.
\]
\end{definition}

$L^2(\mathcal{G})$ is a closed linear subspace of $L^2(\mathcal{F})$. It consists
of exactly those random variables whose values can be determined from the
information available in $\mathcal{G}$---the distinctions that the observational
framework $\mathcal{G}$ makes visible.

\section{Orthogonal Projection as Collapse}

\begin{theorem}[Conditional Expectation as Orthogonal Projection]
For any $X \in L^2(\mathcal{F})$ and any sub-$\sigma$-algebra
$\mathcal{G} \subseteq \mathcal{F}$, the conditional expectation
$\mathbb{E}[X \mid \mathcal{G}]$ is the unique element of $L^2(\mathcal{G})$
that minimizes $\mathbb{E}[(X - Z)^2]$ over all $Z \in L^2(\mathcal{G})$.
Equivalently, $\mathbb{E}[X \mid \mathcal{G}]$ is the orthogonal projection of
$X$ onto $L^2(\mathcal{G})$.
\end{theorem}

\begin{proof}
The closed subspace $L^2(\mathcal{G})$ has a unique nearest point to $X$
by the Hilbert space projection theorem. This nearest point $\hat{X}$ satisfies
the orthogonality condition: $\langle X - \hat{X}, Z \rangle = 0$ for all
$Z \in L^2(\mathcal{G})$. Expanding: $\mathbb{E}[(X - \hat{X})Z] = 0$ for all
$\mathcal{G}$-measurable $Z$. This is precisely Kolmogorov's characterization of
the conditional expectation. By the uniqueness of the projection, $\hat{X} =
\mathbb{E}[X \mid \mathcal{G}]$.
\end{proof}

\begin{definition}[Spherepop Interpretation of Conditional Expectation]
Under the Spherepop identification:
\begin{align*}
\text{History space } \calH &\longleftrightarrow L^2(\mathcal{F}) \\
\text{Observation space } O_c &\longleftrightarrow L^2(\mathcal{G}) \\
\text{Collapse functor } F_c &\longleftrightarrow \mathbb{E}[\cdot \mid \mathcal{G}] \\
\text{Collapse rule } c &\longleftrightarrow \text{Sub-}\sigma\text{-algebra } \mathcal{G} \\
\text{Observational entropy } S_c(o) &\longleftrightarrow \text{Variance of residual}
\end{align*}
\end{definition}

\section{The Orthogonality Condition as State Illusion}

\begin{theorem}[Orthogonality as State Illusion]
The residual $X - \mathbb{E}[X \mid \mathcal{G}]$ is orthogonal to every element
of $L^2(\mathcal{G})$:
\[
\bigl\langle X - \mathbb{E}[X \mid \mathcal{G}],\, Z \bigr\rangle = 0
\quad \text{for all } Z \in L^2(\mathcal{G}).
\]
The discarded component is completely invisible from within the observational
framework that produced the collapse.
\end{theorem}

\begin{proof}
By the projection theorem, $X - \mathbb{E}[X \mid \mathcal{G}]$ is perpendicular
to $L^2(\mathcal{G})$. Expanding: $\mathbb{E}[(X - \mathbb{E}[X\mid\mathcal{G}])Z] =
\mathbb{E}[XZ] - \mathbb{E}[\mathbb{E}[X\mid\mathcal{G}] \cdot Z]$.
By the definition of conditional expectation,
$\mathbb{E}[\mathbb{E}[X\mid\mathcal{G}] \cdot Z] = \mathbb{E}[XZ]$
for all $\mathcal{G}$-measurable $Z$. Therefore the difference is zero.
\end{proof}

\begin{remark}
This theorem is the Hilbert-space version of the state illusion. The observer
with access only to $L^2(\mathcal{G})$ cannot detect the residual component
$X - \mathbb{E}[X\mid\mathcal{G}]$: every measurement they can make is
orthogonal to it, yielding zero correlation. The observer perceives only
$\mathbb{E}[X\mid\mathcal{G}]$ and is unaware that a residual exists. This
is precisely the state illusion: the observer perceives a single observable
state while the underlying history contains information invisible to their
measurement apparatus.
\end{remark}

\chapter{The Algebraic Laws as Geometry}

\section{The Tower Property as Nested Collapse}

The Tower Property of conditional expectation---$\mathbb{E}[\mathbb{E}[X \mid
\mathcal{G}_2] \mid \mathcal{G}_1] = \mathbb{E}[X \mid \mathcal{G}_1]$ when
$\mathcal{G}_1 \subseteq \mathcal{G}_2$---is a geometric theorem about nested
projections.

\begin{theorem}[Tower Property as Nested Projection]
If $\mathcal{G}_1 \subseteq \mathcal{G}_2 \subseteq \mathcal{F}$, then
$L^2(\mathcal{G}_1) \subseteq L^2(\mathcal{G}_2)$, and
\[
\mathbb{E}\bigl[\mathbb{E}[X \mid \mathcal{G}_2] \mid \mathcal{G}_1\bigr]
= \mathbb{E}[X \mid \mathcal{G}_1].
\]
\end{theorem}

\begin{proof}
Since $\mathcal{G}_1 \subseteq \mathcal{G}_2$, every $\mathcal{G}_1$-measurable
function is $\mathcal{G}_2$-measurable. Therefore $L^2(\mathcal{G}_1) \subseteq
L^2(\mathcal{G}_2)$. Let $\hat{X}_2 = \mathbb{E}[X \mid \mathcal{G}_2]$ be
the projection of $X$ onto $L^2(\mathcal{G}_2)$. The residual
$X - \hat{X}_2$ is orthogonal to $L^2(\mathcal{G}_2)$ and hence to the subspace
$L^2(\mathcal{G}_1) \subseteq L^2(\mathcal{G}_2)$. Therefore
$\mathbb{E}[\hat{X}_2 \mid \mathcal{G}_1]$ is the projection of $\hat{X}_2$ onto
$L^2(\mathcal{G}_1)$. But the projection of $X$ onto $L^2(\mathcal{G}_1)$ equals
the projection of $\hat{X}_2 + (X - \hat{X}_2)$ onto $L^2(\mathcal{G}_1)$, and
the residual contributes zero (being orthogonal to $L^2(\mathcal{G}_1)$). Hence
$\mathbb{E}[X \mid \mathcal{G}_1] = \mathbb{E}[\hat{X}_2 \mid \mathcal{G}_1]$.
\end{proof}

\begin{remark}[Spherepop Interpretation]
The Tower Property is the Hilbert-space version of the Spherepop result on
Collapse Functor Composition (Theorem 20.2): applying a coarser collapse after
a finer one is equivalent to applying the coarser collapse directly. In both
settings, the coarser observational framework erases the distinctions preserved
by the finer one. The geometry makes the ``why'' transparent: projecting twice
onto nested subspaces is the same as projecting once onto the smaller subspace.
\end{remark}

\section{The Observational Lattice as Inclusion of Subspaces}

\begin{proposition}[Lattice-Subspace Correspondence]
The lattice of collapse rules ordered by refinement ($c_1 \preceq c_2$ iff
$c_1$ is finer) corresponds to the lattice of closed subspaces of $L^2(\mathcal{F})$
ordered by inclusion: finer rules correspond to larger (more informative)
subspaces.
\[
c_1 \preceq c_2 \quad \iff \quad L^2(\mathcal{G}_{c_1}) \supseteq L^2(\mathcal{G}_{c_2}).
\]
\end{proposition}

\begin{proof}
A finer rule $c_1$ preserves more distinctions, hence its observational subspace
$L^2(\mathcal{G}_{c_1})$ contains more measurable functions---it is a larger
subspace. A coarser rule $c_2$ preserves fewer distinctions---its subspace
$L^2(\mathcal{G}_{c_2})$ is smaller. The inclusion reversal corresponds to the
duality between refinement (more information) and subspace containment (larger
space of representable functions).
\end{proof}

\section{Eve's Law as the Pythagorean Theorem}

The Law of Total Variance (Eve's Law) is not an algebraic identity but a
Pythagorean theorem in $L^2(\mathcal{F})$.

\begin{theorem}[Eve's Law as Pythagorean Theorem]
\[
\mathrm{Var}(X) = \mathbb{E}[\mathrm{Var}(X \mid \mathcal{G})] + \mathrm{Var}(\mathbb{E}[X \mid \mathcal{G}]).
\]
\end{theorem}

\begin{proof}
Center $X$ by writing $\tilde{X} = X - \mathbb{E}[X]$. Decompose:
\[
\tilde{X} = \underbrace{(X - \mathbb{E}[X \mid \mathcal{G}])}_{\text{residual}} +
            \underbrace{(\mathbb{E}[X \mid \mathcal{G}] - \mathbb{E}[X])}_{\text{explained}}.
\]
The residual lies in $L^2(\mathcal{G})^\perp$ (it is the orthogonal complement
of the projection). The explained component lies in $L^2(\mathcal{G})$. These
two vectors are orthogonal. Therefore by the Pythagorean theorem:
\[
\|\tilde{X}\|^2 = \|X - \mathbb{E}[X\mid\mathcal{G}]\|^2 +
                  \|\mathbb{E}[X\mid\mathcal{G}] - \mathbb{E}[X]\|^2.
\]
Taking expectations: $\|\tilde{X}\|^2 = \mathrm{Var}(X)$;
$\|X - \mathbb{E}[X\mid\mathcal{G}]\|^2 = \mathbb{E}[\mathrm{Var}(X\mid\mathcal{G})]$
by the Tower Property; and $\|\mathbb{E}[X\mid\mathcal{G}] - \mathbb{E}[X]\|^2 =
\mathrm{Var}(\mathbb{E}[X\mid\mathcal{G}])$.
\end{proof}

\begin{remark}[Spherepop Interpretation]
Eve's Law decomposes total uncertainty into two orthogonal components:
\begin{itemize}
\item $\mathbb{E}[\mathrm{Var}(X\mid\mathcal{G})]$: \emph{hidden uncertainty}---the
expected residual uncertainty that remains after the observational collapse
$\mathcal{G}$ has occurred. This corresponds to the fibre entropy $S_c(o)$ summed
over observations: the irreducible ambiguity within each equivalence class.
\item $\mathrm{Var}(\mathbb{E}[X\mid\mathcal{G}])$: \emph{visible uncertainty}---the
variance of the observable itself. This corresponds to the variation between
equivalence classes.
\end{itemize}
The total variance $\mathrm{Var}(X)$ is the Pythagorean sum of these two orthogonal
sources. The state illusion hides the first component from the observer who
accesses the history only through $\mathcal{G}$.
\end{remark}

\chapter{Sigma-Algebras as Observational Frameworks}

\section{Information as Geometry}

The $L^2$ framework gives a geometric account of information. An observational
framework $\mathcal{G}$ is not merely a collection of measurable sets; it is a
geometric constraint on what can be known. The subspace $L^2(\mathcal{G})$
consists of exactly those history-valued functions whose values can be determined
given the information in $\mathcal{G}$.

\begin{theorem}[Sigma-Algebra as Information Constraint]
A sub-$\sigma$-algebra $\mathcal{G} \subseteq \mathcal{F}$ corresponds to an
observational framework that makes visible exactly the distinctions in
$L^2(\mathcal{G})$. The following are equivalent:
\begin{enumerate}
\item $X$ is $\mathcal{G}$-measurable.
\item $X \in L^2(\mathcal{G})$.
\item The value of $X$ is completely determined by the observation framework $\mathcal{G}$.
\item $\mathbb{E}[X \mid \mathcal{G}] = X$ (projecting $X$ onto its own subspace leaves it unchanged).
\end{enumerate}
\end{theorem}

\begin{proof}
(1) $\Leftrightarrow$ (2) by definition of $L^2(\mathcal{G})$. (2) $\Leftrightarrow$
(3) because $\mathcal{G}$-measurable functions are exactly those determined by
$\mathcal{G}$-information. (3) $\Leftrightarrow$ (4): if $X \in L^2(\mathcal{G})$,
projecting $X$ onto $L^2(\mathcal{G})$ yields $X$ itself (the projection of a
vector already in a subspace is the vector).
\end{proof}

\section{Collapse Rules and Sub-Sigma-Algebras}

\begin{theorem}[Collapse Rules Are Sub-Sigma-Algebras]
Every Spherepop collapse rule $c : \calH \to O_c$ determines a sub-$\sigma$-algebra
$\mathcal{G}_c$ of the history $\sigma$-algebra $\mathcal{F}$:
\[
\mathcal{G}_c = \{c^{-1}(U) \mid U \subseteq O_c\}.
\]
The collapse functor $F_c$ corresponds to the conditional expectation operator
$\mathbb{E}[\cdot \mid \mathcal{G}_c]$.
\end{theorem}

\begin{proof}
The collection $\{c^{-1}(U) \mid U \subseteq O_c\}$ is closed under
complementation and countable unions (since $c$ is measurable), hence forms a
$\sigma$-algebra. It is a sub-$\sigma$-algebra of $\mathcal{F}$ because every
$c^{-1}(U)$ is a union of fibers of $c$, hence measurable in $\mathcal{F}$.
The conditional expectation $\mathbb{E}[X \mid \mathcal{G}_c]$ then picks out
the component of $X$ that varies only across the equivalence classes of $c$,
which is exactly what the collapse functor $F_c$ does at the combinatorial level.
\end{proof}

\section{The Radon-Nikodym Theorem as Universal Note}

The Radon-Nikodym theorem, which grounds the measure-theoretic definition of
conditional expectation, has a natural interpretation in the Spherepop framework.

\begin{theorem}[Radon-Nikodym as Universal Reconstruction]
Let $\mu$ and $\nu$ be probability measures on $(\Omega, \mathcal{F})$ with
$\nu \ll \mu$ (absolute continuity). Then there exists a unique
$\mathcal{F}$-measurable function $f = d\nu/d\mu$ such that
$\nu(A) = \int_A f \, d\mu$ for all $A \in \mathcal{F}$.
\end{theorem}

\begin{remark}[Spherepop Interpretation]
The Radon-Nikodym derivative $d\nu/d\mu$ is a note in the Spherepop sense: it
preserves the information needed to reconstruct the measure $\nu$ from the
reference measure $\mu$. Without $d\nu/d\mu$, an agent who knows only $\mu$
cannot determine $\nu$. With $d\nu/d\mu$---a measurable function, hence a
Generative Note in the $L^2$ setting---the reconstruction is complete. The
Radon-Nikodym theorem is therefore a universal reconstruction theorem: it
guarantees the existence of the minimal note needed to recover one measure from
another, whenever such recovery is possible.
\end{remark}

\section{Variance as Continuous Observational Entropy}

\begin{theorem}[Variance as Continuous Fibre Entropy]
The conditional variance $\mathrm{Var}(X \mid \mathcal{G})$ is the continuous
analogue of the discrete observational entropy $S_c(o) = \log |c^{-1}(o)|$:
\[
\mathrm{Var}(X \mid \mathcal{G}) = \mathbb{E}\bigl[(X - \mathbb{E}[X \mid \mathcal{G}])^2 \mid \mathcal{G}\bigr]
\]
measures the expected squared deviation of $X$ from its projection onto
$\mathcal{G}$, which is the expected information content of the residual
component invisible to the observational framework.
\end{theorem}

\begin{proof}
The residual $X - \mathbb{E}[X\mid\mathcal{G}]$ is the component of $X$
that the framework $\mathcal{G}$ cannot see. Its conditional variance
$\mathbb{E}[(X - \mathbb{E}[X\mid\mathcal{G}])^2 \mid \mathcal{G}]$ measures
the expected squared magnitude of this invisible component, given the visible
component. This is the continuous analogue of counting the histories in a
fibre: instead of $\log |c^{-1}(o)|$ discrete histories, we have a continuous
measure of the spread of histories within the fibre.
\end{proof}

\begin{remark}[Unification]
This theorem completes the bridge between the discrete Spherepop collapse theory
and the continuous $L^2$ conditional expectation theory. The dictionary is:

\begin{center}
\begin{tabular}{ll}
\toprule
Discrete Spherepop & Continuous $L^2$ \\
\midrule
History space $\calH$ & $L^2(\mathcal{F})$ \\
Collapse rule $c$ & Sub-$\sigma$-algebra $\mathcal{G}$ \\
Collapse functor $F_c$ & Conditional expectation $\mathbb{E}[\cdot\mid\mathcal{G}]$ \\
Observational entropy $S_c(o)$ & Conditional variance $\mathrm{Var}(X\mid\mathcal{G})$ \\
Fibre $c^{-1}(o)$ & Fiber over $\mathcal{G}$-atom \\
State illusion & Orthogonality of residual \\
Note (anti-collapse) & Radon-Nikodym derivative \\
Tower Property & Nested projection theorem \\
Eve's Law & Pythagorean theorem in $L^2$ \\
\bottomrule
\end{tabular}
\end{center}
\end{remark}

\part{Open Problems and Future Directions}

\chapter{The Inverse Problem and Observational Completeness}

\section{When Does Observational Equivalence Imply History Equivalence?}

The deepest open question in the Spherepop framework is the converse of the
collapse functor: when does $c(H_1) = c(H_2)$ imply $H_1 = H_2$?

\begin{definition}[Observational Completeness]
A collapse rule $c$ is \emph{observationally complete} if it is injective:
$c(H_1) = c(H_2) \Rightarrow H_1 = H_2$.
\end{definition}

\begin{definition}[Collapse Kernel]
The \emph{kernel} of a collapse rule $c$ is
\[
\ker(c) = \{(H_1, H_2) \in \calH \times \calH \mid c(H_1) = c(H_2)\}.
\]
Observational completeness holds iff $\ker(c) = \{(H, H) \mid H \in \calH\}$.
\end{definition}

\begin{proposition}[Completeness Results for Canonical Rules]
The Identity rule $c_I$ is observationally complete. The LastWrite rule $c_{LW}$
is not: $[pop(x), pop(y)]$ and $[pop(y), pop(x), pop(y)]$ are
$c_{LW}$-equivalent but distinct. The Accumulate rule $c_{\mathrm{acc}}$ is not:
$[pop(x), pop(y)]$ and $[pop(y), pop(x)]$ have identical accumulate-states.
\end{proposition}

\begin{definition}[Inverse-Problem Fibre]
For $o \in O_c$, the \emph{fibre} of $o$ is
$c^{-1}(o) = \{H \in \calH \mid c(H) = o\}$.
\end{definition}

Characterising fibres for general rules, and determining when a fibre contains a
unique history, is the primary open mathematical problem in the framework.

\section{Open Problems}

\subsection{Distributed Spherepop Runtimes}

The current World is single-threaded. A distributed Spherepop runtime would
coordinate multiple Worlds, each with their own histories, under a global
consistency constraint. The consistency constraint must be expressed in terms of
history equivalence under a shared collapse rule---not value consistency. The
event sourcing and CRDT literatures provide partial models; the Spherepop-specific
challenge is the certified-refusal structure across distributed agents.

\subsection{Collapse Functor Composition}

The four canonical collapse rules are special cases of a more general family of
functors $F_c : \Hist \to \Obs_c$. A theory of functor composition would
characterise which combinations of rules are valid---when $F_{c_2} \circ F_{c_1}$
is itself a valid collapse functor---and would classify the algebraic structure
of the space of collapse rules.

\subsection{Dependent Process Types in Full Generality}

The current Process$(T_1 \rightsquigarrow T_2)$ type is first-order. A dependent
process type $\Pi(x : T_1).\, \mathrm{Process}(x, f(x))$ would allow the type
of the output to depend on the specific value consumed. This is the natural
extension of dependent type theory into the process calculus setting and would
enable precise typing of programs that branch on their inputs.

\subsection{Admissibility Lattice Completeness}

The current type system uses a Boolean admissibility structure. The full theory
requires a graded admissibility lattice with continuous admissibility levels,
connecting to the RSVP xylomorphic criterion $\lambda < 1$ and the admissibility
geometry program. Completeness theorems for the graded type system would
characterize exactly which trajectories are admissible at each level.

\subsection{Proof-Carrying Refusal in Full Generality}

The current $\mathsf{RefusalReason}$ is a vocabulary. A complete proof-carrying
refusal system would make $\mathsf{RefusalReason}$ a proof term: evidence that
can be type-checked independently of the computation that produced it. This
connects to the theory of certified refusal and would enable mechanized
verification of inadmissibility claims.

\chapter{The Coda: Implementation as Philosophical Argument}

The implementation did not merely realize the theory; it exposed which distinctions
the theory was still hiding.

The central claim of this monograph is not that Spherepop is a good programming
language, though we believe it is an interesting one. The central claim is that
taking a philosophical position seriously enough to implement it is a form of
philosophical argument.

Four implementation discoveries were not predicted by the initial framework. They
were forced by the attempt to mechanize it.

The original $\Refuse(\mathsf{Symbol})$ gestured at documented inadmissibility.
It took implementation pressure to reveal that a symbol alone documents nothing:
the reason is the document.

The original notion of observable state gestured at the quotient structure. It
took \texttt{World::observe} to force the question: observable under what rule?

The original type checker gestured at world-relative knowledge. It took the
free-variable conflict to force the question: which world are we assuming?

The original correctness criterion gestured at semantic preservation. It took the
compiler test to force the question: preservation of what, exactly?

Each discovery reveals a distinction the philosophical framework was gesturing at
without making precise. Only in the attempt to mechanize---to force the philosophy
to compile---did the gestures become definitions. This is the deepest lesson of
ontology-driven language design. Philosophy can identify the right questions. Only
implementation can force the right precision.

\medskip

Spherepop is the algebra of continuations.

Notes are the physical realization of that algebra.

Civilization is the attempt, always incomplete, always under threat, to build a
note infrastructure adequate to the continuation needs of the species that requires
it.

\newpage


%% ============================================================
%% PART X: CALCULUS OF CONSTRUCTIONS
%% ============================================================
\part{The Calculus of Constructions}

\chapter{Why Spherepop Needs the Calculus of Constructions}

\section{The Limits of Simple Types}

The type system developed in Part~IV is powerful enough to certify admissibility,
document refusals, and seal observations. But it is limited in one fundamental
respect: types cannot depend on values. The type of the output of a function cannot
depend on what the input actually is. This limitation becomes acute in several
situations that arise naturally in Spherepop.

Consider a function that reads a file and returns a proof that the file was
successfully read. The output type should mention the specific file that was read,
not just the abstract type of files. Or consider a GC pass that returns a history
with only live events: the output type should certify that specific inadmissibility
certificates from the input are preserved. Simple types cannot express either of
these conditions.

The Calculus of Constructions (CoC), introduced by Coquand and Huet~\cite{coquand88},
is the natural remedy. It is a typed lambda calculus in which types may depend on
terms, terms may be types, and the distinction between levels collapses into a
single universe hierarchy. In CoC, the type of a returned value can carry the full
specification of what was computed.

\section{The Pure Type System Perspective}

The Calculus of Constructions belongs to the family of \emph{pure type systems}
(PTSs), which unify the handling of terms and types through a single syntactic
category.

\begin{definition}[Sorts and Pure Type System]
A \emph{pure type system} is determined by a triple $(\mathcal{S}, \mathcal{A},
\mathcal{R})$ where $\mathcal{S}$ is a set of \emph{sorts}, $\mathcal{A}$ is a
set of \emph{axioms} $s_1 : s_2$, and $\mathcal{R}$ is a set of \emph{rules}
$(s_1, s_2, s_3)$ governing when $\Pi$-types can be formed.
\end{definition}

\begin{definition}[Calculus of Constructions]
The Calculus of Constructions is the PTS with:
\begin{align*}
\mathcal{S} &= \{\ast, \square\} \\
\mathcal{A} &= \{\ast : \square\} \\
\mathcal{R} &= \{(\ast, \ast, \ast),\; (\ast, \square, \square),\;
                 (\square, \ast, \ast),\; (\square, \square, \square)\}
\end{align*}
where $\ast$ is the sort of ordinary types and $\square$ is the sort of kinds
(types of types).
\end{definition}

The four rules in $\mathcal{R}$ permit:
\begin{enumerate}
\item $(\ast, \ast, \ast)$: functions from terms to terms (ordinary functions).
\item $(\ast, \square, \square)$: functions from terms to types (dependent types).
\item $(\square, \ast, \ast)$: functions from types to terms (polymorphism).
\item $(\square, \square, \square)$: functions from types to types (type operators).
\end{enumerate}

The full power of CoC is that all four are available simultaneously. This is what
makes it capable of expressing the entire type-theoretic apparatus needed for
Spherepop's admissibility certificates.

\section{The Spherepop-CoC Correspondence}

The Spherepop type system is naturally embedded into CoC.

\begin{definition}[CoC Embedding of Spherepop Types]
The embedding $\llbracket \cdot \rrbracket : \mathsf{SphType} \to \mathsf{CoCTerm}$
is defined by:
\begin{align*}
\llbracket \mathsf{Unit} \rrbracket &= \top_{\mathsf{CoC}} \\
\llbracket \mathsf{Never} \rrbracket &= \bot_{\mathsf{CoC}} \\
\llbracket \Adm(T) \rrbracket &= \Sigma(h : \mathsf{History}).\, \mathsf{Admissible}(h, \llbracket T \rrbracket) \\
\llbracket \Rfsd(T, r) \rrbracket &= \Sigma(h : \mathsf{History}).\, \mathsf{Refused}(h, \llbracket T \rrbracket, r) \\
\llbracket \Col(T, c) \rrbracket &= \Sigma(h : \mathsf{History}).\, \mathsf{Collapsed}(h, \llbracket T \rrbracket, c) \\
\llbracket \Pi(x : T_1).T_2 \rrbracket &= \Pi(x : \llbracket T_1 \rrbracket).\llbracket T_2 \rrbracket \\
\llbracket \Proc(T_1 \rightsquigarrow T_2) \rrbracket &=
  \Pi(h_1 : \llbracket T_1 \rrbracket).\, \Sigma(h_2 : \llbracket T_2 \rrbracket).\,
  \mathsf{BindDep}(h_1, h_2)
\end{align*}
\end{definition}

The $\Sigma$-types in the embedding are crucial: they pair the value with a
proof that the history supporting it has the required structure. The admissibility
certificate is not merely a type annotation but a proof term that can be
independently verified.

\begin{theorem}[Embedding Soundness]
If $\Gamma \vdash t : T$ in the Spherepop type system, then
$\llbracket \Gamma \rrbracket \vdash \llbracket t \rrbracket : \llbracket T \rrbracket$
in CoC.
\end{theorem}

\begin{proof}
By structural induction on the Spherepop typing derivation. Each typing rule maps
to a valid CoC derivation:

$[\mathrm{T\text{-}POP}]$: The embedding sends $\mathsf{pop}(t)$ to a pair
$\langle h, \pi \rangle$ where $\pi$ is a proof of $\mathsf{Admissible}(h, \llbracket T \rrbracket)$.
This is a valid $\Sigma$-introduction in CoC.

$[\mathrm{T\text{-}REFUSE}]$: The embedding sends $\mathsf{refuse}(t, r)$ to a
pair $\langle h, \pi_r \rangle$ where $\pi_r$ proves $\mathsf{Refused}(h, \llbracket T \rrbracket, r)$.
Again a valid $\Sigma$-introduction.

$[\mathrm{T\text{-}COLLAPSE}]$: Collapse requires a proof of admissibility as
input and produces a proof of the collapsed type. The CoC derivation uses
$\Sigma$-elimination on the admissibility proof followed by $\Sigma$-introduction
for the collapsed type.

The remaining cases follow by induction.
\end{proof}

\chapter{Implementing CoC for Spherepop}

\section{The Universe Hierarchy}

A full implementation of CoC for Spherepop requires a universe hierarchy to avoid
Russell-type paradoxes. We adopt a predicative hierarchy.

\begin{definition}[Universe Hierarchy]
Define sorts $\mathsf{Type}_0, \mathsf{Type}_1, \mathsf{Type}_2, \ldots$ with
axioms $\mathsf{Type}_i : \mathsf{Type}_{i+1}$ and cumulative coercions
$\mathsf{Type}_i \hookrightarrow \mathsf{Type}_{i+1}$.
\end{definition}

In practice, most Spherepop typing occurs at $\mathsf{Type}_0$ (the sort of
ordinary program types). The proof-carrying refusal reasons inhabit
$\mathsf{Type}_1$ (propositions about $\mathsf{Type}_0$ terms). The collapse
rules inhabit $\mathsf{Type}_1$ as functions from histories to observable states.

\begin{lstlisting}[style=rust, caption={Universe levels in the Spherepop type checker}]
pub enum Universe {
    Type(usize),   // Type_0, Type_1, ...
    Prop,          // proof-irrelevant propositions
}

pub enum Term {
    Sort(Universe),
    Var(Symbol),
    App(Box<Term>, Box<Term>),
    Lam(Symbol, Box<Term>, Box<Term>),  // λx:A.b
    Pi(Symbol, Box<Term>, Box<Term>),   // Πx:A.B
    Sigma(Symbol, Box<Term>, Box<Term>),// Σx:A.B
    Pair(Box<Term>, Box<Term>),         // (a, b)
    Fst(Box<Term>),                     // π₁
    Snd(Box<Term>),                     // π₂
    // Spherepop primitives
    Pop(Box<Term>),
    Refuse(Box<Term>, RefusalReason),
    Bind(Box<Term>, Box<Term>),
    Collapse(Box<Term>, CollapseRule),
    // Admissibility universe
    Admissible(Box<Term>),
    Refused(Box<Term>, RefusalReason),
    Collapsed(Box<Term>, CollapseRule),
}
\end{lstlisting}

\section{The Type Checker}

A CoC type checker for Spherepop must implement bidirectional type checking:
type inference for introduction forms and type checking for elimination forms.

\begin{definition}[Bidirectional Typing]
Bidirectional typing uses two judgements:
\begin{itemize}
\item $\Gamma \vdash t \Rightarrow T$ (\emph{synthesis}): the type $T$ of $t$
is inferred from $t$ alone.
\item $\Gamma \vdash t \Leftarrow T$ (\emph{checking}): $t$ is checked against
the expected type $T$.
\end{itemize}
\end{definition}

\begin{lstlisting}[style=rust, caption={Bidirectional type checker core}]
pub fn synthesize(ctx: &Context, term: &Term)
    -> Result<Term, TypeError>
{
    match term {
        Term::Var(x) => ctx.lookup(x),
        Term::App(f, a) => {
            let f_ty = synthesize(ctx, f)?;
            match whnf(&f_ty) {
                Term::Pi(x, a_ty, b_ty) => {
                    check(ctx, a, &a_ty)?;
                    Ok(subst(b_ty, x, a))
                }
                _ => Err(TypeError::NotAFunction)
            }
        }
        Term::Fst(p) => {
            let p_ty = synthesize(ctx, p)?;
            match whnf(&p_ty) {
                Term::Sigma(_, a, _) => Ok(*a),
                _ => Err(TypeError::NotAPair)
            }
        }
        Term::Pop(t) => {
            let t_ty = synthesize(ctx, t)?;
            Ok(Term::Admissible(Box::new(t_ty)))
        }
        Term::Refuse(t, r) => {
            let t_ty = synthesize(ctx, t)?;
            Ok(Term::Refused(Box::new(t_ty), r.clone()))
        }
        Term::Collapse(t, c) => {
            let t_ty = synthesize(ctx, t)?;
            match whnf(&t_ty) {
                Term::Admissible(inner) =>
                    Ok(Term::Collapsed(inner, c.clone())),
                _ => Err(TypeError::CollapseRequiresAdmissible)
            }
        }
        _ => Err(TypeError::CannotSynthesize)
    }
}

pub fn check(ctx: &Context, term: &Term, ty: &Term)
    -> Result<(), TypeError>
{
    match (term, whnf(ty)) {
        (Term::Lam(x, _, b), Term::Pi(y, a, b_ty)) => {
            let ctx2 = ctx.extend(x, &a);
            check(&ctx2, b, &subst(&b_ty, y, &Term::Var(x)))
        }
        (Term::Pair(a, b), Term::Sigma(x, a_ty, b_ty)) => {
            check(ctx, a, &a_ty)?;
            let a_val = eval(a);
            check(ctx, b, &subst(&b_ty, x, &a_val))
        }
        _ => {
            let t_ty = synthesize(ctx, term)?;
            if definitionally_equal(&t_ty, ty) { Ok(()) }
            else { Err(TypeError::TypeMismatch) }
        }
    }
}
\end{lstlisting}

\section{Definitional Equality and Normalization}

CoC type checking requires deciding definitional equality: two types are
definitionally equal if they reduce to the same normal form.

\begin{definition}[Weak Head Normal Form]
A term is in \emph{weak head normal form} (WHNF) if it is not a beta-redex at
the head: it is a lambda, a Pi, a Sigma, a sort, a variable, or an application
whose function is in WHNF.
\end{definition}

\begin{lstlisting}[style=rust, caption={WHNF reduction for Spherepop-CoC}]
pub fn whnf(term: &Term) -> Term {
    match term {
        Term::App(f, a) => {
            match whnf(f) {
                Term::Lam(x, _, b) => whnf(&subst(&b, &x, a)),
                f_whnf => Term::App(Box::new(f_whnf),
                                    a.clone())
            }
        }
        Term::Fst(p) => {
            match whnf(p) {
                Term::Pair(a, _) => whnf(&a),
                p_whnf => Term::Fst(Box::new(p_whnf))
            }
        }
        Term::Snd(p) => {
            match whnf(p) {
                Term::Pair(_, b) => whnf(&b),
                p_whnf => Term::Snd(Box::new(p_whnf))
            }
        }
        // Spherepop reductions
        Term::Pop(t) => Term::Pop(Box::new(whnf(t))),
        Term::Collapse(t, c) => {
            Term::Collapse(Box::new(whnf(t)), c.clone())
        }
        t => t.clone()
    }
}
\end{lstlisting}

\begin{theorem}[Normalization of Spherepop-CoC]
The Spherepop-CoC type system is strongly normalizing: every well-typed
reduction sequence terminates.
\end{theorem}

\begin{proof}[Proof sketch]
The Calculus of Constructions is strongly normalizing by the reducibility
candidates argument of Girard~\cite{girard89}. The Spherepop extensions Pop,
Refuse, Bind, and Collapse are all introduction forms that do not introduce
new beta-redexes. They reduce only by unwrapping their arguments to WHNF.
The combined system inherits strong normalization from CoC by a standard
extension argument: the new rules add no new reduction paths that could create
infinite sequences, since the underlying lambda calculus reductions are already
terminating.
\end{proof}

\section{Identity Types and History Equality}

CoC extended with identity types (Martin-L\"of type theory~\cite{martinlof84})
gives Spherepop a native notion of proof-level history equality.

\begin{definition}[Identity Type]
For any type $A$ and terms $a, b : A$, the \emph{identity type} $a =_A b$ is a
type whose inhabitants are proofs that $a$ and $b$ are definitionally equal.
The sole constructor is $\mathsf{refl}_a : a =_A a$.
\end{definition}

\begin{definition}[History Identity Type]
For histories $H_1, H_2 : \mathsf{History}$, the type $H_1 =_{\mathsf{History}} H_2$
is inhabited precisely when the histories are event-for-event identical. This is
the type-level statement of Replay Equivalence: a term of type
$I(p).\mathsf{history} =_{\mathsf{History}} V(C(p)).\mathsf{history}$
is a formal proof that the compiler is correct.
\end{definition}

\begin{lstlisting}[style=rust, caption={History identity in the type system}]
pub fn verify_replay_equivalence(
    prog: &Term,
    interp_world: &World,
    compiled_world: &World,
) -> Result<(), ProofError> {
    // Produce a proof of history identity
    if interp_world.history == compiled_world.history {
        Ok(())  // refl witnesses the identity
    } else {
        // Construct a counterexample witness
        let diff = history_diff(
            &interp_world.history,
            &compiled_world.history
        );
        Err(ProofError::HistoryMismatch(diff))
    }
}
\end{lstlisting}

\section{The Curry-Howard Correspondence for Spherepop}

The Curry-Howard correspondence---that proofs are programs and propositions are
types---takes a specific form in Spherepop.

\begin{proposition}[Spherepop Curry-Howard Correspondence]
The following identifications hold:
\begin{center}
\begin{tabular}{lll}
\toprule
Logical Notion & Program Notion & Spherepop Realization \\
\midrule
Proposition & Type & $T : \mathsf{Type}_0$ \\
Proof & Term & $t : T$ \\
Implication & Function type & $\Pi(x:A).B$ \\
Conjunction & Pair type & $\Sigma(x:A).B$ \\
Admissibility & Certificate type & $\Adm(T)$ \\
Refutation & Refusal type & $\Rfsd(T, r)$ \\
Observation & Collapsed type & $\Col(T, c)$ \\
History identity & Replay equivalence & $H_1 =_\mathcal{H} H_2$ \\
\bottomrule
\end{tabular}
\end{center}
\end{proposition}

%% ============================================================
%% PART XI: THE BNF GRAMMAR AS SUPERIOR NOTATION
%% ============================================================
\part{The BNF Grammar as Universal Notation}

\chapter{Why Grammar Beats Circuits and Code}

\section{The Problem of Notation}

Every computational formalism needs a notation: a way of writing down the objects
it studies. The choice of notation is not merely aesthetic. It determines what
can be expressed without ambiguity, what can be processed mechanically, what can
be read without specialized tools, and what can be implemented independently of
any particular hardware or software substrate.

Three dominant notations for computation exist: programming language syntax,
circuit diagrams, and formal grammars. The thesis of this chapter is that
Backus-Naur Form (BNF) grammar, or its extensions, is the superior notation for
specifying computational behavior, and that Spherepop's use of BNF is not a
convenience but a theoretical necessity.

\section{A Brief History of BNF}

Backus-Naur Form was introduced by Peter Naur in the ALGOL 60 report~\cite{naur60}.
Its purpose was to give a precise, implementation-independent specification of a
programming language's syntax. Before BNF, programming language grammars were
described in English prose: ambiguous, difficult to parse mechanically, and
dependent on shared human conventions.

BNF replaces prose with a minimal formal apparatus: nonterminals (enclosed in
angle brackets), terminals (written as themselves), a production arrow ($::=$),
and alternation ($\mid$). From these four elements, any context-free language can
be specified.

\begin{definition}[BNF Grammar]
A \emph{Backus-Naur Form grammar} is a tuple $G = (N, T, P, S)$ where:
\begin{itemize}
\item $N$ is a finite set of \emph{nonterminals}.
\item $T$ is a finite set of \emph{terminals} (disjoint from $N$).
\item $P$ is a finite set of \emph{productions} of the form $A ::= \alpha$
where $A \in N$ and $\alpha \in (N \cup T)^*$.
\item $S \in N$ is the \emph{start symbol}.
\end{itemize}
\end{definition}

Extended BNF (EBNF) adds optional elements ($[\cdot]$), repetition ($\{\cdot\}$),
and grouping, but adds no expressive power beyond context-free languages.

\section{BNF Versus Circuit Diagrams}

Circuit diagrams are the standard notation for hardware computation. They have
significant advantages: they are visually intuitive, they make data flow explicit,
and they can be directly compiled to hardware. But they have several critical
disadvantages relative to BNF.

\subsection{Implementation Dependence}

A circuit diagram is inherently tied to a specific level of abstraction. A gate
diagram for a full adder specifies AND, OR, and NOT gates in a specific topology.
This specification cannot be directly used at a different level of abstraction
without redrawing. A BNF grammar for the same operation specifies the structure
independently of implementation:

\begin{lstlisting}[style=rust, caption={BNF for binary addition (implementation-independent)}]
<bit>     ::= "0" | "1"
<addend>  ::= <bit> | <bit> <addend>
<sum>     ::= <addend> "+" <addend>
<carry>   ::= <bit>
<result>  ::= <carry> <sum>
\end{lstlisting}

This grammar specifies the structure of binary addition without committing to
gates, lookup tables, carry-lookahead circuits, or any other implementation.
Any implementation that accepts the grammar's language is correct by definition.

\subsection{Human Readability}

Circuit diagrams require specialized visual parsing. They cannot be read aloud,
stored in plain text, or processed by standard text tools. A BNF grammar is
written in ASCII (or Unicode), can be stored in a version control repository,
diffed, searched, and processed by any text processor. It is read left-to-right,
top-to-bottom, in the same direction as natural language.

\begin{theorem}[BNF Readability Advantage]
For any grammar $G$ specifying a language $L$, the BNF representation of $G$
requires only the characters of $N \cup T$ plus the BNF metasymbols
$\{::=, |, \langle, \rangle\}$. The circuit representation of the same computation
requires a graphical medium that cannot be stored or transmitted as plain text.
\end{theorem}

\begin{proof}
BNF is defined over a finite alphabet. Its productions are strings over that
alphabet. A circuit diagram, by contrast, represents spatial relationships (the
topology of the circuit) that cannot be serialized as a string without loss of
the spatial information that makes the diagram legible. The best-known lossless
serialization of a circuit diagram is a netlist, which is itself a grammar-like
notation---a point in favor of the grammar approach.
\end{proof}

\subsection{Compositional Structure}

BNF grammars compose naturally. If $G_1$ specifies language $L_1$ and $G_2$
specifies $L_2$, the grammar for $L_1 \cdot L_2$ (concatenation) is formed by
adding a production $S ::= S_1\, S_2$ and combining the production sets. This
compositional structure is the foundation of modular language design.

\begin{proposition}[Grammar Compositionality]
If $G_1 = (N_1, T_1, P_1, S_1)$ and $G_2 = (N_2, T_2, P_2, S_2)$ with
$N_1 \cap N_2 = \varnothing$, then:
\begin{align*}
G_1 \cdot G_2 &= (N_1 \cup N_2 \cup \{S\},\, T_1 \cup T_2,\,
                  P_1 \cup P_2 \cup \{S ::= S_1\, S_2\},\, S) \\
G_1 \mid G_2 &= (N_1 \cup N_2 \cup \{S\},\, T_1 \cup T_2,\,
                  P_1 \cup P_2 \cup \{S ::= S_1 \mid S_2\},\, S)
\end{align*}
These operations correspond to sequential composition and choice in Spherepop.
\end{proposition}

\section{BNF Versus Programming Language Syntax}

BNF is more fundamental than any particular programming language syntax because
it is what programming languages are defined in. A programming language's syntax
is a grammar; BNF is the metalanguage for grammars.

But the comparison is not merely definitional. BNF grammars have concrete
advantages as computational specifications over program text.

\subsection{Ambiguity is Detectable}

A BNF grammar can be checked for ambiguity (multiple parse trees for the same
string). An English prose description of a language cannot. A program in any
conventional language can be syntactically unambiguous but semantically ambiguous
in ways that BNF would surface if the semantics were given as a grammar.

\subsection{Separation of Concerns}

BNF cleanly separates the \emph{shape} of a program (its grammar) from its
\emph{meaning} (its semantics). This separation is enforced by the formalism:
a BNF grammar says nothing about what programs mean, only about what strings
are valid programs. This is not a limitation but a virtue: it allows the same
grammar to be given multiple semantics (operational, denotational, axiomatic)
without changing the specification of valid programs.

\subsection{Formal Verification Target}

A BNF grammar is a mathematical object that can serve as the specification for
formal verification. A Spherepop program can be verified to be within the language
of a grammar by a parser; verification at the semantic level then proceeds over
the parse tree. This two-stage verification is not possible when the language is
specified in prose or implicitly in a reference implementation.

\section{The Spherepop BNF as Computational Universal}

The Spherepop BNF grammar from Part~III is not merely a syntactic convenience.
It is a specification from which an implementation in any language, on any
substrate, can be mechanically derived.

\begin{definition}[Spherepop Core Grammar]
The Spherepop core language is given by the following BNF:
\begin{align*}
\langle\mathit{term}\rangle &::= \langle\mathit{var}\rangle
  \mid \langle\mathit{lambda}\rangle
  \mid \langle\mathit{app}\rangle
  \mid \langle\mathit{let}\rangle
  \mid \langle\mathit{pop}\rangle
  \mid \langle\mathit{refuse}\rangle
  \mid \langle\mathit{bind}\rangle
  \mid \langle\mathit{collapse}\rangle
  \mid \langle\mathit{seq}\rangle \\
\langle\mathit{var}\rangle &::= \langle\mathit{ident}\rangle \\
\langle\mathit{lambda}\rangle &::= (\lambda \mid \mathsf{fn})\;
  \langle\mathit{ident}\rangle\; [:\langle\mathit{type}\rangle]\;
  (\to \mid .)\; \langle\mathit{term}\rangle \\
\langle\mathit{app}\rangle &::= \langle\mathit{term}\rangle\;
  \langle\mathit{term}\rangle \\
\langle\mathit{let}\rangle &::= \mathsf{let}\; \langle\mathit{ident}\rangle\;
  = \langle\mathit{term}\rangle\; \mathsf{in}\; \langle\mathit{term}\rangle \\
\langle\mathit{pop}\rangle &::= \mathsf{pop}\; \langle\mathit{term}\rangle \\
\langle\mathit{refuse}\rangle &::= \mathsf{refuse}\; \langle\mathit{term}\rangle\;
  [\langle\mathit{reason}\rangle] \\
\langle\mathit{bind}\rangle &::= \mathsf{bind}\; \langle\mathit{term}\rangle\;
  \langle\mathit{term}\rangle \\
\langle\mathit{collapse}\rangle &::= \mathsf{collapse}\;
  \langle\mathit{term}\rangle\; [\langle\mathit{rule}\rangle] \\
\langle\mathit{seq}\rangle &::= \mathsf{seq}\;
  [\langle\mathit{term}\rangle\; \{;\langle\mathit{term}\rangle\}^*] \\
\langle\mathit{reason}\rangle &::= \mathsf{violation}(\langle\mathit{ident}\rangle)
  \mid \mathsf{explicit}(\langle\mathit{ident}\rangle) \\
\langle\mathit{rule}\rangle &::= \mathsf{id} \mid \mathsf{last\_write}
  \mid \mathsf{accumulate} \mid \mathsf{proj}(\langle\mathit{ident}\rangle) \\
\langle\mathit{type}\rangle &::= \mathsf{Unit} \mid \mathsf{Never}
  \mid \langle\mathit{type\_var}\rangle
  \mid \Adm(\langle\mathit{type}\rangle)
  \mid \langle\mathit{type}\rangle \to \langle\mathit{type}\rangle
\end{align*}
\end{definition}

\begin{theorem}[Grammar Determines Implementation]
Given the Spherepop core BNF, a parser, type checker, and evaluator are
mechanically derivable for any target platform.
\end{theorem}

\begin{proof}
The grammar is context-free and therefore parseable by any LALR(1) or PEG
parser generator. The type checker follows from the typing rules (which are
themselves a kind of grammar over typing contexts). The evaluator follows from
the operational semantics (which are also a grammar of world-state transitions).
Each of these derivations is mechanical and platform-independent: the same BNF
yields implementations in Rust, Python, Haskell, JavaScript, or any other host
language without modification to the grammar.
\end{proof}

\subsection{Grammar as Interface Contract}

When the Spherepop grammar is fixed, it constitutes an interface contract between:
(1) program authors, who write terms in the grammar's language;
(2) parsers, which accept exactly the grammar's language;
(3) type checkers, which verify admissibility of parsed terms;
(4) evaluators, which execute well-typed terms;
(5) compilers, which translate well-typed terms to Event IR.

Any implementation satisfying this contract is a correct Spherepop implementation.
The grammar is the single source of truth for what the language is, independent
of any particular reference implementation.

\section{BNF as Anti-Collapse Artifact}

The relationship between BNF grammars and the note theory of this monograph is
not accidental.

A BNF grammar is a Refusal Note of the highest generativity. It specifies which
strings are admissible (members of the language) and which are inadmissible. It
refuses the collapse of the language into any particular implementation. And it
generates the full continuation space of all programs in the language: any string
in the language is a reachable continuation, and the grammar specifies exactly
the set of reachable continuations.

\begin{theorem}[BNF Grammars as Generative Refusal Notes]
A BNF grammar $G$ is simultaneously:
\begin{enumerate}
\item A \emph{Generative Note}: it creates the continuation space $L(G)$ of all
programs in the language, which did not exist before the grammar was specified.
\item A \emph{Refusal Note}: it refuses all strings not in $L(G)$ as
inadmissible, with the implicit reason $\mathsf{ParseError}$.
\item A \emph{Coordination Note}: it couples the behaviors of all
implementations by binding them to the same accepted and rejected strings.
\item A \emph{Capture Note}: it preserves the structure of valid programs across
implementations, allowing programs written for one implementation to be run
on another.
\end{enumerate}
\end{theorem}

\begin{proof}
(1) The language $L(G)$ is the set of all strings derivable from $S$ in $G$.
Before $G$ exists, this set is undefined; after $G$ exists, it is precisely
specified. New programs in $L(G)$ are continuations that $G$ makes reachable.

(2) A string $w \notin L(G)$ has no parse tree under $G$. The parser refuses
it with a parse error. This is a Refuse operation with reason
$\mathsf{ParseError}$ applied to the inadmissible string.

(3) Any two implementations $I_1$ and $I_2$ of $G$ accept the same strings and
reject the same strings. They are bound by the grammar to the same
admissibility decisions. This is a Bind coupling their trajectory spaces.

(4) A program $p \in L(G)$ has an implementation-independent parse tree. The
parse tree captures the structure of $p$ independently of what evaluator
processes it. This is a Capture Note for the program's syntactic structure.
\end{proof}

\chapter{Grammar, Circuits, and the Efficiency Hierarchy}

\section{Measuring Expressiveness Per Symbol}

We want a precise way to compare notations for expressing computational structure.
Let the \emph{expressiveness per symbol} of a notation $\mathcal{N}$ for
specification $S$ be the ratio of the information content of $S$ to the number
of symbols required to express $S$ in $\mathcal{N}$.

\begin{definition}[Specification Complexity]
The \emph{specification complexity} $K_\mathcal{N}(S)$ of specification $S$ in
notation $\mathcal{N}$ is the minimum number of symbols in $\mathcal{N}$ required
to express $S$ unambiguously. The \emph{relative expressiveness} of $\mathcal{N}_1$
over $\mathcal{N}_2$ for $S$ is
\[
E(S, \mathcal{N}_1, \mathcal{N}_2) = \frac{K_{\mathcal{N}_2}(S)}{K_{\mathcal{N}_1}(S)}.
\]
$E > 1$ means $\mathcal{N}_1$ is more expressive per symbol.
\end{definition}

\begin{theorem}[BNF Specification Efficiency]
For context-free languages, BNF is asymptotically more symbol-efficient than
circuit diagrams and at least as efficient as any other notation that is
implementation-independent and mechanically parseable.
\end{theorem}

\begin{proof}[Proof sketch]
A circuit diagram for a function $f : \{0,1\}^n \to \{0,1\}^m$ specifies $f$
by encoding its topology as a directed acyclic graph of gates. The minimum size
of such a diagram is $\Omega(C(f))$ where $C(f)$ is the circuit complexity of
$f$. A BNF grammar for the same function, specified as a language recognizer,
can be written in $O(\log n + \log m + K(f))$ productions where $K(f)$ is the
Kolmogorov complexity of $f$'s description. For functions with short
descriptions but high circuit complexity (such as sorting networks), BNF can be
exponentially more compact.
\end{proof}

\section{The Readability Spectrum}

Different notations occupy different positions on the spectrum from machine-efficient
to human-efficient.

\begin{center}
\begin{tabular}{llll}
\toprule
Notation & Human Readable & Machine Parseable & Impl.-Independent \\
\midrule
Machine code & No & Yes & No \\
Assembly & Partial & Yes & No \\
Circuit diagram & Yes (visual) & Partial & No \\
C/Rust program & Yes & Yes & Partial \\
Lambda calculus & Yes & Yes & Yes \\
BNF grammar & Yes & Yes & Yes \\
Natural language & Yes & No & Yes \\
\bottomrule
\end{tabular}
\end{center}

BNF occupies the unique position of being simultaneously human-readable,
mechanically parseable, and fully implementation-independent. No other common
notation achieves all three.

\section{Grammar as the Missing Level of Abstraction}

Programming languages, circuit diagrams, and prose all operate at specific
levels of abstraction. Grammar operates at the level of \emph{specification}:
it describes the structure of all possible instances of a computational artifact
without committing to any particular instance.

\begin{definition}[Abstraction Level Hierarchy]
The abstraction levels of computational description form a hierarchy:
\begin{enumerate}
\item \emph{Physical}: transistor layouts, voltages, timing.
\item \emph{Circuit}: gate-level topology.
\item \emph{Machine}: instruction set architecture.
\item \emph{Language}: program syntax and semantics.
\item \emph{Grammar}: specification of the language itself.
\item \emph{Meta-grammar}: specification of the notation for grammars.
\end{enumerate}
\end{definition}

\begin{proposition}[Grammar as Implementation Firewall]
A grammar at level $n$ of the hierarchy is implementation-independent with
respect to all levels below $n$. A Spherepop BNF grammar is therefore
independent of all language implementations, all machine instruction sets, all
circuit designs, and all physical substrates.
\end{proposition}

This independence is exactly what makes BNF appropriate as the specification
language for Spherepop. The Spherepop grammar specifies what computations the
language can express without committing to how they are executed. The execution
details---Rust runtime, WASM compilation, hardware acceleration---are below the
abstraction level at which the grammar operates.

\section{Literate Grammar and Documentation}

Knuth's literate programming~\cite{knuth84} proposed that programs should be
written as literature: prose explanations interleaved with code. The Spherepop
approach inverts this slightly: grammar should be written as literature, with
prose explanations interleaved with BNF productions.

\begin{definition}[Literate Grammar]
A \emph{literate grammar} is a BNF grammar augmented with:
\begin{enumerate}
\item Prose explanations of the intent of each production.
\item Examples of derivations in the grammar.
\item Semantic notes indicating what well-formed strings mean.
\item Transformation rules mapping productions to type-checking and evaluation rules.
\end{enumerate}
\end{definition}

The Spherepop BNF in this monograph is a literate grammar in this sense: each
production is accompanied by the operational semantics rule it corresponds to
and the type-checking rule it enables. This makes the grammar simultaneously
a specification, a tutorial, and a formal verification target.

%% ============================================================
%% PART XII: SUPPLEMENTARY DERIVATIONS
%% ============================================================
\part{Supplementary Derivations}

\chapter{Derivation Compendium}

This chapter collects all supplementary derivations, filling the mathematical
gaps indicated in the preceding parts. Each derivation is labeled by the chapter
it supports.

\section{Chapter 1: Notes Increase Reachability}

For each continuation $c \in \calC$, define the \emph{reachability gain} of note
$N$ by
\[
G_A(N, c, t) = P_A(c, t \mid N) - P_A(c, t).
\]
Then $N$ is a note precisely when $\exists c \in \calC,\; G_A(N, c, t) > 0$.

\begin{proposition}[Nontrivial Notes Are Reachability-Increasing]
If $N$ is a note and all non-note effects are nonnegative, then the total
reachability gain
$G_A(N, t) = \sum_{c \in \calC} G_A(N, c, t)$
is strictly positive.
\end{proposition}

\begin{proof}
By definition, there exists at least one $c_0$ such that $G_A(N, c_0, t) > 0$.
If all other terms satisfy $G_A(N, c, t) \geq 0$, then the total sum contains
one strictly positive term and no negative terms. Hence $G_A(N, t) > 0$.
\end{proof}

\section{Chapter 2: State as Quotient}

\begin{proposition}[State Is a Quotient of History]
Every surjective state projection $\sigma : \calH \to S$ factors uniquely through
the quotient map $q : \calH \to \calH/{\sim_\sigma}$.
\end{proposition}

\begin{proof}
Define $\bar\sigma([H]) = \sigma(H)$. This is well-defined because
$H_1 \sim_\sigma H_2$ implies $\sigma(H_1) = \sigma(H_2)$. Surjectivity makes
$\bar\sigma$ onto; injectivity follows because distinct classes have distinct
$\sigma$-values. Thus $S \cong \calH/{\sim_\sigma}$.
\end{proof}

\section{Chapter 3: Operator Minimality}

\begin{theorem}[Operator Minimality]
No one of $\Pop, \Refuse, \Bind, \Collapse$ is definable from the other three
while preserving its effect on history, option space, admissibility, and
observation.
\end{theorem}

\begin{proof}
$\Pop$ uniquely decreases $\Omega$; the other three leave $\Omega$ unchanged.
$\Refuse$ uniquely changes the admissibility set $\calA(H)$ without decreasing
$\Omega$; $\Pop$ changes $\Omega$, while $\Bind$ and $\Collapse$ do not record
inadmissibility. $\Bind$ uniquely introduces a dependency relation between two
elements without consuming or refusing either; no combination of the other
operators produces a pure dependency edge. $\Collapse$ uniquely maps a history
through an observational rule into a quotient space; the other operators extend
history internally but produce no observation. Hence all four are irreducible.
\end{proof}

\section{Chapter 4: Free Category}

\begin{theorem}[Histories Form the Free Category on the Event Graph]
Let $G$ be the directed graph whose edges are primitive events. Then $\Hist$ is
the free category generated by $G$.
\end{theorem}

\begin{proof}
Objects are option spaces; morphisms are finite paths. Given any category
$\mathcal{D}$ and graph morphism $F : G \to U(\mathcal{D})$, the unique functor
$\bar{F} : \Hist \to \mathcal{D}$ extending $F$ sends $[e_1, \ldots, e_n]$ to
$F(e_n) \circ \cdots \circ F(e_1)$. This is the universal property of the free
category.
\end{proof}

\section{Chapter 5: Small-Step Determinacy}

\begin{proposition}[Determinacy Under Fixed Rule Selection]
If the next operator and its arguments are fixed, Spherepop small-step reduction
is deterministic.
\end{proposition}

\begin{proof}
Each transition rule has a unique conclusion for fixed inputs. No fixed operator
application can reduce to two distinct worlds.
\end{proof}

\section{Chapter 6: Collapse Entropy Monotonicity}

\begin{theorem}[Collapse Increases Observational Entropy Under Coarsening]
If $c_2$ is coarser than $c_1$, then every $c_2$-observation has entropy at
least as large as the entropy of any $c_1$-observation contained inside it.
\end{theorem}

\begin{proof}
Coarsening means each $c_1$-class is contained in some $c_2$-class. The fibre
of the coarser observation contains the fibre of the finer observation, so
$|c_2^{-1}(o_2)| \geq |c_1^{-1}(o_1)|$ for contained observations. Taking
logarithms gives the entropy inequality.
\end{proof}

\section{Chapter 7: Type Soundness}

\begin{theorem}[Spherepop Type Soundness]
If $\vdash t : T$ and $t \to^* t'$, then either $t'$ is a value of type $T$,
or there exists $t''$ such that $t' \to t''$.
\end{theorem}

\begin{proof}
By repeated application of Preservation and Progress. Preservation gives
$\vdash t' : T$ after any finite reduction sequence. Progress implies $t'$ is
either a value or can take another step.
\end{proof}

\section{Chapter 8: Lattice Collapse Threshold}

\begin{proposition}[Monotonicity of Collapse Permission]
If $\ell_1 \leq \ell_2$ and $\Gamma \vdash t : \Adm_{\ell_2}(T)$, then every
collapse permitted at level $\ell_1$ is also permitted at level $\ell_2$.
\end{proposition}

\begin{proof}
A higher admissibility level satisfies at least the constraints of a lower level.
Since $\ell_1 \leq \ell_2$, any rule whose threshold is met by $\ell_1$ is also
met by $\ell_2$.
\end{proof}

\section{Note Composition Theory}

Define note composition by
\[
\calC(N_1 \otimes N_2) = \calC(N_1) \cup \calC(N_2) \cup \calC_{\mathrm{bind}}(N_1, N_2),
\]
where $\calC_{\mathrm{bind}}(N_1, N_2)$ is the set of continuations reachable
only because both notes are jointly available.

\begin{theorem}[Coordination Surplus]
If $\calC_{\mathrm{bind}}(N_1, N_2) \neq \varnothing$, then
\[
V_C(N_1 \otimes N_2) > V_C(N_1) + V_C(N_2) - |\calC(N_1) \cap \calC(N_2)|.
\]
\end{theorem}

\begin{proof}
The continuation set of the composite contains the union of the individual sets
plus at least one additional bound continuation. The ordinary union has size
$V_C(N_1) + V_C(N_2) - |\calC(N_1) \cap \calC(N_2)|$. Adding the nonempty
bound set strictly increases this.
\end{proof}

\begin{theorem}[Note Entropy]
Define the \emph{note entropy} of $N$ relative to agent $A$ at time $t$ as
\[
H_A(N, t) = -\sum_{c \in \calC} P_A(c, t \mid N) \log P_A(c, t \mid N).
\]
A note $N'$ is more \emph{focused} than $N$ if $H_A(N', t) < H_A(N, t)$: it
concentrates reachability on a smaller set of continuations. A note $N'$ is more
\emph{generative} than $N$ if $V_C(N') > V_C(N)$.
\end{theorem}

\begin{remark}
Focus and generativity trade off: collapse notes are highly focused (low entropy)
while generative notes are high-entropy. The note taxonomy orders by this tradeoff.
\end{remark}

\section{Note Class Derivations}

\begin{proposition}[Capture Adequacy]
A capture note $N$ is adequate for task family $T$ iff for every continuation
$c \in T$ there exists recoverable data $d \subseteq N$ such that
$P_A(c, t \mid d) > P_A(c, t)$.
\end{proposition}

\begin{theorem}[Prospective Notes Are Deferred Pops]
Every prospective note determines a deferred commitment operator. At time $t$
the continuation is not yet executed (no immediate Pop occurs); at time $t'$ the
agent follows the note, selecting the future continuation via a Pop.
\end{theorem}

\begin{proof}
A prospective note binds present state $A_t$ to future continuation $c_{t'}$.
The selection at $t'$ is a Pop of $c_{t'}$ from the available alternatives.
The note stores a commitment whose execution is deferred. \end{proof}

\begin{theorem}[Repair Notes Minimize Reconstruction Cost]
A repair note $N$ is optimal in family $\mathcal{F}$ when
$C_r(c \mid N) = \min_{M \in \mathcal{F}} C_r(c \mid M)$.
\end{theorem}

\begin{proof}
Repair value is $C_r(c) - C_r(c \mid N)$. Since $C_r(c)$ is fixed for the
continuation, maximizing repair value is equivalent to minimizing $C_r(c \mid N)$.
\end{proof}

\begin{proposition}[Checksums as Pure Refusal Notes]
A checksum $\chi(h)$ refuses informational collapse by distinguishing corrupted
histories from admissible ones. If $\chi(h') \neq \chi(h)$, then $h'$ is refused
as a valid continuation of $h$. The checksum is a pure refusal note: it stores
only the boundary between admissible and inadmissible continuations.
\end{proposition}

\begin{theorem}[Generativity Is Strict Reachability Expansion]
A note $N$ is generative iff $\exists c : P_A(c,t) = 0$ and $P_A(c,t \mid N) > 0$.
\end{theorem}

\begin{proof}
If such $c$ exists, $N$ creates access to a previously unreachable trajectory.
Conversely, generativity requires opening at least one previously unavailable
trajectory, which is exactly this condition.
\end{proof}

\begin{theorem}[Collapse Notes Trade Detail for Access]
Let $\calN$ be a finite note collection. A collapse note $K$ partitioning $\calN$
into $m$ classes reduces search cost from $O(|\calN|)$ to $O(m + \max_i |\calN_i|)$.
\end{theorem}

\begin{proof}
Without organizing structure, worst-case search examines every note. With a
partition of $m$ classes, the agent selects one class ($m$ steps) then searches
that class ($\max_i |\calN_i|$ steps worst case).
\end{proof}

\section{Compilation and GC Derivations}

\begin{theorem}[Event IR Is History-Preserving]
For every primitive source event $e$, compilation followed by VM execution appends
the same event to history as interpretation.
\end{theorem}

\begin{proof}
By case analysis: $\Pop$ compiles to \texttt{IrEvent::Pop}, which the VM appends
as $\Pop$; $\Refuse(r)$ compiles to \texttt{IrEvent::Refuse\{r\}}, appended as
$\Refuse(r)$; analogously for $\Bind$ and $\Collapse(c)$.
\end{proof}

\begin{theorem}[GC Is Safe Exactly When It Preserves Live Fibres]
A GC operation $g : \calH \to \calH$ is safe iff for every live continuation $c$,
$P_A(c, t \mid g(H)) = P_A(c, t \mid H)$.
\end{theorem}

\begin{proof}
If the equality holds for all live continuations, GC has removed no history needed
for any reachable future---hence safe. Conversely, safety by definition preserves
all history required for live continuations.
\end{proof}

\section{Error Correction Derivation}

\begin{theorem}[Certified Refusal Prevents Silent Error Absorption]
If every refusal carries a proof term $\pi$, then no error can be erased without
either being handled or remaining visible in the output type.
\end{theorem}

\begin{proof}
A refusal has type $\Rfsd(T, \pi)$. The typing rules allow this type to be handled
by a recovery term or propagated. There is no rule converting $\Rfsd(T, \pi)$ into
$T$ while discarding $\pi$. Hence the error cannot be silently absorbed.
\end{proof}

\section{Category Theory Derivations}

\begin{theorem}[Collapse as Reflection]
If $\Obs_c$ is the full subcategory of $\Hist$ consisting of histories saturated
under $\sim_c$, then the quotient functor $F_c : \Hist \to \Obs_c$ is left
adjoint to the inclusion $I_c : \Obs_c \hookrightarrow \Hist$.
\end{theorem}

\begin{proof}
For each history $H$ and saturated observable $O$, every morphism $H \to I_c(O)$
constant on $\sim_c$-classes factors uniquely through $F_c(H)$. This gives the
natural bijection $\Obs_c(F_c(H), O) \cong \Hist(H, I_c(O))$, establishing
$F_c \dashv I_c$.
\end{proof}

\begin{corollary}[Perfect Observation Requires History]
An observation map is perfectly history-faithful iff it is isomorphic to the
identity observation on $\calH$.
\end{corollary}

\begin{proof}
Perfect history-faithfulness requires injectivity. A bijective observation
establishes an isomorphism between histories and observations, hence is isomorphic
to the identity.
\end{proof}

\section{Sheaf Derivations}

\begin{theorem}[State Illusion as Failure of Unique Gluing]
Given a cover of observations $\{U_i\}$, state illusion occurs exactly when
compatible local sections $s_i \in \calH(U_i)$ admit more than one global
section.
\end{theorem}

\begin{proof}
Compatible local sections agree on overlaps. If more than one global history
restricts to the same local observations, the observations do not determine a
unique history---precisely the state illusion.
\end{proof}

\section{CLIO Derivation}

\begin{theorem}[Spherepop Observational Entropy Is Discrete CLIO Entropy]
The CLIO fibre entropy $S_\pi(o) = \log \mathrm{Vol}(\pi^{-1}(o))$ reduces in
the discrete setting to $S_c(o) = \log |c^{-1}(o)|$.
\end{theorem}

\begin{proof}
In the discrete setting, volume is counting measure. Therefore
$\mathrm{Vol}(\pi^{-1}(o)) = |\pi^{-1}(o)| = |c^{-1}(o)|$.
\end{proof}

\section{Active Geodesic Derivation}

\begin{proposition}[Histories Are Discrete Geodesics]
Let $L(e_i)$ be the local cost of event $e_i$ and define the action
$\mathcal{S}(H) = \sum_i L(e_i)$. A Spherepop execution is a discrete geodesic
when it minimizes $\mathcal{S}(H)$ among admissible histories with the same
boundary conditions.
\end{proposition}

\begin{proof}
In discrete variational mechanics, a geodesic minimizes action between fixed
endpoints. A Spherepop history is a path in event space; restricting to admissible
histories and minimizing the additive cost functional gives the discrete geodesic
condition.
\end{proof}

\section{MEM\textbar8 Locality Derivation}

\begin{theorem}[Locality as Continuation Preservation]
If related memory components are physically or logically local, then the expected
cost of reconstructing their shared continuation is lower than under random
placement.
\end{theorem}

\begin{proof}
Let $d(x, y)$ be access distance between memory components. Reconstruction cost
is monotone in aggregate distance among components required by a continuation.
Local placement lowers expected pairwise distance relative to random placement,
hence lowers expected reconstruction cost.
\end{proof}

\section{Civilizational Collapse Derivation}

\begin{theorem}[Reachability Loss Under Archive Destruction]
Let $c \in R(Z)$ be reachable only through some $N \in \calN$. Destroying $\calN$
strictly decreases collective reachability.
\end{theorem}

\begin{proof}
Since $c$ is reachable only through $N$, removing $N$ removes $c$ from $R(Z)$.
The post-destruction reachability is a proper subset of the original.
\end{proof}

\section{Inverse Problem Derivation}

\begin{proposition}[Kernel Triviality Characterizes Complete Observation]
A collapse rule $c$ is observationally complete iff
$\ker(c) = \Delta_\calH = \{(H, H) : H \in \calH\}$.
\end{proposition}

\begin{proof}
Observational completeness means $c(H_1) = c(H_2) \Rightarrow H_1 = H_2$, so
the kernel contains only diagonal pairs. Conversely, a diagonal kernel gives
injectivity.
\end{proof}

\section{Final Unification Theorem}

\begin{theorem}[Notes as Reachability Artifacts]
\label{thm:unification}
An artifact $N$ is a note iff it is a reachability artifact: there exists an
agent $A$, a continuation $c$, and a future time $t$ such that
$P_A(c, t \mid N) > P_A(c, t)$.
Spherepop programs are notes because they preserve, restrict, bind, refuse, and
collapse continuations in history space.
\end{theorem}

\begin{proof}
The first claim is the definition of a note under the continuation account. For
the second: every Spherepop program constructs a history from primitive events.
$\Pop$ selects continuations; $\Refuse$ preserves inadmissibility distinctions;
$\Bind$ couples trajectories; $\Collapse$ makes histories observable under rules.
A Spherepop program is therefore an artifact that changes which continuations
remain reachable to future agents. Hence it is a note.
\end{proof}

\begin{corollary}[Universal Reachability Artifact Thesis]
Programs, proofs, archives, maps, protocols, libraries, institutions, languages,
and civilizations are all special cases of the same abstract object: a reachability
artifact. They differ in the medium through which they preserve continuations and
in the scale at which they operate. The abstract structure is identical in every
case.
\end{corollary}

\begin{proof}
By Theorem~\ref{thm:unification}, any artifact that increases the probability of
traversing at least one continuation is a note. Each of the listed artifacts does
this:

A \emph{program} preserves a computational continuation.
A \emph{proof} preserves an inferential continuation.
An \emph{archive} preserves historical continuations against decay.
A \emph{map} creates navigational continuations.
A \emph{protocol} binds communicative continuations between agents.
A \emph{library} is a repository of continuations across domains.
An \emph{institution} binds social and coordination continuations across
generations.
A \emph{language} is a system of semantic continuations shared by a community.
A \emph{civilization} is the totality of all of the above, organized into a
distributed hierarchy of note infrastructure.

In every case, the note increases $P_A(c, t)$ for some continuation $c$.
\end{proof}

\begin{thebibliography}{99}

\bibitem{barendregt84}
Barendregt, H.~P. \emph{The Lambda Calculus: Its Syntax and Semantics}.
North-Holland, 1984.

\bibitem{borgman07}
Borgman, C.~L. \emph{Scholarship in the Digital Age}. MIT Press, 2007.

\bibitem{clark98}
Clark, A. and Chalmers, D. The extended mind. \emph{Analysis} 58, no.~1 (1998):
7--19.

\bibitem{eisenstein80}
Eisenstein, E.~L. \emph{The Printing Press as an Agent of Change}. Cambridge
University Press, 1980.

\bibitem{feldman58}
Feldman, J. On the absolute continuity of Gaussian measures. \emph{Zeitschrift
f\"ur Wahrscheinlichkeitstheorie} 1 (1958).

\bibitem{felleisen87}
Felleisen, M. and Friedman, D.~P. A calculus for assignments in higher-order
languages. \emph{Proceedings of the 14th ACM SIGACT-SIGPLAN Symposium on
Principles of Programming Languages}, 1987.


\bibitem{chenoweth25}
Chenoweth, C., Chenoweth, A.
MEM\textbar8: A Wave-Based Cognitive Architecture for Multimodal Memory Integration
and Consciousness Simulation. Zenodo, 2025.
\texttt{https://doi.org/10.5281/zenodo.16436298}

\bibitem{barendregt84}
Barendregt, H.~P. \emph{The Lambda Calculus: Its Syntax and Semantics}.
North-Holland, 1984.

\bibitem{borgman07}
Borgman, C.~L. \emph{Scholarship in the Digital Age}. MIT Press, 2007.

\bibitem{coquand88}
Coquand, T. and Huet, G. The calculus of constructions.
\emph{Information and Computation} 76 (1988): 95--120.

\bibitem{clark98}
Clark, A. and Chalmers, D. The extended mind. \emph{Analysis} 58, no.~1 (1998):
7--19.

\bibitem{eisenstein80}
Eisenstein, E.~L. \emph{The Printing Press as an Agent of Change}. Cambridge
University Press, 1980.

\bibitem{feldman58}
Feldman, J. On the absolute continuity of Gaussian measures. \emph{Zeitschrift
f\"ur Wahrscheinlichkeitstheorie} 1 (1958).

\bibitem{felleisen87}
Felleisen, M. and Friedman, D.~P. A calculus for assignments in higher-order
languages. \emph{Proceedings of the 14th ACM SIGACT-SIGPLAN Symposium on
Principles of Programming Languages}, 1987.

\bibitem{fowler05}
Fowler, M. Event sourcing. \texttt{martinfowler.com}, 2005.

\bibitem{girard89}
Girard, J.-Y., Lafont, Y., and Taylor, P. \emph{Proofs and Types}. Cambridge
University Press, 1989.

\bibitem{goody77}
Goody, J. \emph{The Domestication of the Savage Mind}. Cambridge University
Press, 1977.

\bibitem{hajek58}
H\'ajek, J. On a property of normal distributions of an arbitrary stochastic
process. \emph{Czechoslovak Mathematical Journal} 8 (1958).

\bibitem{hoare85}
Hoare, C.~A.~R. \emph{Communicating Sequential Processes}. Prentice Hall, 1985.

\bibitem{hofmann98}
Hofmann, M. and Streicher, T. The groupoid interpretation of type theory.
In \emph{Twenty-five Years of Constructive Type Theory}. Oxford University Press,
1998.

\bibitem{hutchins95}
Hutchins, E. \emph{Cognition in the Wild}. MIT Press, 1995.

\bibitem{knuth84}
Knuth, D.~E. Literate programming. \emph{The Computer Journal} 27, no.~2
(1984): 97--111.

\bibitem{maclane71}
Mac Lane, S. \emph{Categories for the Working Mathematician}. Springer, 1971.

\bibitem{martinlof84}
Martin-L\"of, P. \emph{Intuitionistic Type Theory}. Bibliopolis, 1984.

\bibitem{milner80}
Milner, R. \emph{A Calculus of Communicating Systems}. Springer, 1980.

\bibitem{naur60}
Naur, P. (ed.). Report on the algorithmic language ALGOL 60.
\emph{Communications of the ACM} 3, no.~5 (1960): 299--314.
(Original introduction of Backus--Naur Form.)

\bibitem{nordstrom90}
Nordstr\"om, B., Petersson, K., and Smith, J. \emph{Programming in
Martin-L\"of's Type Theory}. Oxford University Press, 1990.

\bibitem{ong82}
Ong, W.~J. \emph{Orality and Literacy: The Technologizing of the Word}.
Methuen, 1982.

\bibitem{pierce02}
Pierce, B.~C. \emph{Types and Programming Languages}. MIT Press, 2002.

\bibitem{reiter78}
Reiter, R. On closed world data bases. In \emph{Logic and Data Bases}. Plenum
Press, 1978.

\bibitem{santoro24}
Santoro, M., Waghmare, S., and Panaretos, V.~M. Kernel embeddings of functional
data and mutual singularity. Preprint, 2024.

\bibitem{shapiro11}
Shapiro, M., Pregui\c{c}a, N., Baquero, C., and Zawirski, M. Conflict-free
replicated data types. INRIA Technical Report, 2011.

\bibitem{strachey67}
Strachey, C. Fundamental concepts in programming languages. \emph{Higher-Order
and Symbolic Computation} 13 (2000): 11--49. Originally delivered 1967.

\bibitem{tulving72}
Tulving, E. Episodic and semantic memory. In \emph{Organization of Memory}.
Academic Press, 1972.

\bibitem{univalent13}
The Univalent Foundations Program. \emph{Homotopy Type Theory: Univalent
Foundations of Mathematics}. Institute for Advanced Study, 2013.

\bibitem{wright94}
Wright, A.~K. and Felleisen, M. A syntactic approach to type soundness.
\emph{Information and Computation} 115, no.~1 (1994): 38--94.

\end{thebibliography}


\end{document}
